% !TEX program = lualatex
% ---------------------------------------------------------------------------------------------
%  Readback note 26: Good forms and rows: equal Hecke traces on unitary groups that differ at the real places
%
%  Build:  python3 tools/rb.py note <slug>
%
%  One unit of a Lean artifact, rewritten as ordinary mathematics. The prose is hand-written. Lean names,
%  line ranges and counts come from data.tex; quoted Lean statements are copied from the artifact by the build.
%
%  Rules of the text: standard terminology; every symbol is explained where it first appears and listed again
%  in the notation at the end; no letter has two meanings; a proof is complete or says what it takes for granted;
%  Lean names appear only in the LEAN lines and in the concordance.
%
%  The signature rule (the build enforces it): a function, relation or predicate that the artifact defines is
%  written only through its notation macro, which the build makes from the spec (`notation`) with one slot for
%  each explicit argument of the Lean definition. No argument is fixed and dropped, also not "for brevity" and
%  not for one section. Every argument and every hypothesis of a stated declaration stands inside its
%  statement, or in a Setting that the statement cites by number; only implicit and instance arguments (the
%  ambient objects) may stand in the conventions.
% ---------------------------------------------------------------------------------------------
\documentclass[11pt,a4paper]{article}
\ifdefined\pdfvariable\pdfvariable suppressoptionalinfo \numexpr 1+2+4+8\relax\fi
\usepackage[a4paper,hmargin=35mm,top=22mm,bottom=27mm,footskip=12mm,nohead]{geometry}

\usepackage{fontspec}
\directlua{
  luaotfload.add_fallback("gfallback", {"STIXTwoMath-Regular.otf:mode=node;", "DejaVuSans.ttf:mode=node;"})
  luaotfload.add_fallback("gmonofallback", {"DejaVuSansMono.ttf:mode=node;", "STIXTwoMath-Regular.otf:mode=node;"})
}
\setmainfont{NewCM10}[Extension=.otf,UprightFont=*-Book,ItalicFont=*-BookItalic,
  BoldFont=*-Bold,BoldItalicFont=*-BoldItalic,SmallCapsFont=*-Book,SmallCapsFeatures={Letters=SmallCaps},
  RawFeature={fallback=gfallback}]
\setsansfont{NewCMSans10}[Extension=.otf,UprightFont=*-Book,ItalicFont=*-BookOblique,
  BoldFont=*-Bold,BoldItalicFont=*-BoldOblique,RawFeature={fallback=gfallback}]
\setmonofont{JetBrainsMono}[Extension=.otf,UprightFont=*-Regular,ItalicFont=*-Italic,
  BoldFont=*-SemiBold,BoldItalicFont=*-SemiBoldItalic,Scale=MatchLowercase,
  RawFeature={-calt;-liga;fallback=gmonofallback}]
\usepackage{amsmath}
\usepackage{amsthm}
\usepackage[math-style=TeX]{unicode-math}
\setmathfont{NewCMMath-Book.otf}
\emergencystretch=2.5em\relax

\usepackage{xcolor}
\usepackage{tikz}
\usetikzlibrary{calc,arrows.meta,positioning,backgrounds}
\usepackage[protrusion=true,expansion=true]{microtype}
\usepackage{enumitem}
\usepackage{needspace}
\usepackage{booktabs}
\usepackage{array}
\usepackage{longtable}
\usepackage{fancyhdr}
\usepackage{lastpage}
\usepackage{multicol}
\usepackage[hidelinks,pdfusetitle]{hyperref}

\definecolor{ink}{HTML}{0F172A}
\definecolor{inksoft}{HTML}{475569}
\definecolor{inkfaint}{HTML}{94A3B8}
\definecolor{rulegrey}{HTML}{CBD5E1}
\definecolor{arith}{HTML}{C2410C}
\definecolor{est}{HTML}{7C3AED}
\definecolor{ana}{HTML}{0891B2}
\definecolor{ok}{HTML}{047857}
\definecolor{leancol}{HTML}{3730A3}

\title{Good forms and rows: equal Hecke traces on unitary groups that differ at the real places}
\author{grh-compact}
\date{}

% ------------------------------------------------------------------------------------- macros
\newcommand{\leanfont}{\ttfamily\fontsize{8.6}{10.5}\selectfont\color{leancol}}
\newcommand{\lean}[1]{\csname LN@#1\endcsname}
\newcommand{\fillin}[1]{{\color{arith}\itshape[#1]}}  % a place of the template that is not written yet; the build reports every one that is left
\newcommand{\leanquote}[1]{\csname LQ@#1\endcsname}
% generated by rb.py note; do not edit
\newcommand{\ConcDecls}{70}
\newcommand{\Pin}{72fead10}
\newcommand{\PinText}{that of commit {\leanfont 72fead10}}
\newcommand{\ArtifactName}{grh-compact}
\expandafter\newcommand\csname LN@existscommonheckerallis\endcsname{\href{https://github.com/lsernau/grh-lean/blob/72fead102f663d5c7a407964a13dc0cf9d977207/grh-compact/src/New/Trace/GoodFormTransfer.lean\#L27-L47}{{\leanfont Data.\allowbreak{}exists\_common\_hecke\_rallis}}}
\expandafter\newcommand\csname LQ@existscommonheckerallis\endcsname{\leanquoteline{0}{lemma exists\_common\_hecke\_rallis ϑ j \{a a'\} \{B B' : \_ → \_ → ℂ\} (h : ∀ z y,}\leanquoteline{2}{HasSum (a * D.rallisTerm ϑ j V z y ·) (B z y))}\leanquoteline{2}{(h' : ∀ z y, HasSum (a' * D.rallisTerm ϑ j \_ z y ·) (B' z y)) : ∃ p : D.heckeAlgebra,}\leanquoteline{2}{(∀ z y, B z y = a * inner ℂ z (p.1 V y)) ∧ ∀ z y, B' z y = a' * inner ℂ z (p.1 V' y)}}
\expandafter\newcommand\csname LN@traceheckeOnGoodeq\endcsname{\href{https://github.com/lsernau/grh-lean/blob/72fead102f663d5c7a407964a13dc0cf9d977207/grh-compact/src/New/Trace/GoodFormTransfer.lean\#L56-L89}{{\leanfont Data.\allowbreak{}trace\_heckeOnGood\_eq}}}
\expandafter\newcommand\csname LQ@traceheckeOnGoodeq\endcsname{\leanquoteline{0}{lemma trace\_heckeOnGood\_eq p : trace ℂ \_ (D.heckeOnGood V p) = trace ℂ \_ (D.heckeOnGood V' p)}}
\expandafter\newcommand\csname LN@detheckeOnGoodeq\endcsname{\href{https://github.com/lsernau/grh-lean/blob/72fead102f663d5c7a407964a13dc0cf9d977207/grh-compact/src/New/Trace/GoodFormTransfer.lean\#L94-L98}{{\leanfont Data.\allowbreak{}det\_heckeOnGood\_eq}}}
\expandafter\newcommand\csname LQ@detheckeOnGoodeq\endcsname{\leanquoteline{0}{lemma det\_heckeOnGood\_eq p : (D.heckeOnGood V p).det = (D.heckeOnGood V' p).det}}
\expandafter\newcommand\csname LN@goodFormsinfeigenParteqbot\endcsname{\href{https://github.com/lsernau/grh-lean/blob/72fead102f663d5c7a407964a13dc0cf9d977207/grh-compact/src/New/Trace/GoodFormTransfer.lean\#L103-L110}{{\leanfont Data.\allowbreak{}goodForms\_inf\_eigenPart\_eq\_bot}}}
\expandafter\newcommand\csname LN@finrankrowseqfinrankgoodForms\endcsname{\href{https://github.com/lsernau/grh-lean/blob/72fead102f663d5c7a407964a13dc0cf9d977207/grh-compact/src/New/Trace/GoodFormTransfer.lean\#L116-L118}{{\leanfont Data.\allowbreak{}finrank\_rows\_eq\_finrank\_goodForms}}}
\expandafter\newcommand\csname LQ@finrankrowseqfinrankgoodForms\endcsname{\leanquoteline{0}{lemma finrank\_rows\_eq\_finrank\_goodForms (j : D.LiftIndex) : D.rowDim =}\leanquoteline{4}{finrank ℂ (D.goodForms (D.spaceW j))}}
\expandafter\newcommand\csname LN@besselIndexEquiv\endcsname{\href{https://github.com/lsernau/grh-lean/blob/72fead102f663d5c7a407964a13dc0cf9d977207/grh-compact/src/New/Trace/GoodIdempotent.lean\#L16-L19}{{\leanfont Shape.\allowbreak{}besselIndexEquiv}}}
\expandafter\newcommand\csname LN@meqtwomulvarrhoaddtwo\endcsname{\href{https://github.com/lsernau/grh-lean/blob/72fead102f663d5c7a407964a13dc0cf9d977207/grh-compact/src/New/Trace/GoodIdempotent.lean\#L27-L27}{{\leanfont Data.\allowbreak{}m\_eq\_two\_mul\_varrho\_add\_two}}}
\expandafter\newcommand\csname LN@levelNormmemandnotMemofmemQ\endcsname{\href{https://github.com/lsernau/grh-lean/blob/72fead102f663d5c7a407964a13dc0cf9d977207/grh-compact/src/New/Trace/GoodIdempotent.lean\#L34-L36}{{\leanfont Data.\allowbreak{}levelNorm\_mem\_and\_notMem\_of\_mem\_Q}}}
\expandafter\newcommand\csname LN@levelPlaces\endcsname{\href{https://github.com/lsernau/grh-lean/blob/72fead102f663d5c7a407964a13dc0cf9d977207/grh-compact/src/New/Trace/GoodIdempotent.lean\#L40-L42}{{\leanfont Data.\allowbreak{}levelPlaces}}}
\expandafter\newcommand\csname LN@decl\endcsname{\href{https://github.com/lsernau/grh-lean/blob/72fead102f663d5c7a407964a13dc0cf9d977207/grh-compact/src/New/Trace/GoodIdempotent.lean\#L46-L46}{{\leanfont Data.\allowbreak{}𝔪}}}
\expandafter\newcommand\csname LN@nonemptymodulusplaces\endcsname{\href{https://github.com/lsernau/grh-lean/blob/72fead102f663d5c7a407964a13dc0cf9d977207/grh-compact/src/New/Trace/GoodIdempotent.lean\#L48-L50}{{\leanfont Data.\allowbreak{}nonempty\_modulus\_places}}}
\expandafter\newcommand\csname LN@chosenIdele\endcsname{\href{https://github.com/lsernau/grh-lean/blob/72fead102f663d5c7a407964a13dc0cf9d977207/grh-compact/src/New/Trace/GoodIdempotent.lean\#L54-L55}{{\leanfont Data.\allowbreak{}chosenIdele}}}
\expandafter\newcommand\csname LN@ideleA\endcsname{\href{https://github.com/lsernau/grh-lean/blob/72fead102f663d5c7a407964a13dc0cf9d977207/grh-compact/src/New/Trace/GoodIdempotent.lean\#L59-L59}{{\leanfont Data.\allowbreak{}ideleA}}}
\expandafter\newcommand\csname LN@piIdele\endcsname{\href{https://github.com/lsernau/grh-lean/blob/72fead102f663d5c7a407964a13dc0cf9d977207/grh-compact/src/New/Trace/GoodIdempotent.lean\#L62-L62}{{\leanfont Data.\allowbreak{}piIdele}}}
\expandafter\newcommand\csname LN@piIdelespec\endcsname{\href{https://github.com/lsernau/grh-lean/blob/72fead102f663d5c7a407964a13dc0cf9d977207/grh-compact/src/New/Trace/GoodIdempotent.lean\#L64-L66}{{\leanfont Data.\allowbreak{}piIdele\_spec}}}
\expandafter\newcommand\csname LN@latticeScale\endcsname{\href{https://github.com/lsernau/grh-lean/blob/72fead102f663d5c7a407964a13dc0cf9d977207/grh-compact/src/New/Trace/GoodIdempotent.lean\#L70-L70}{{\leanfont Data.\allowbreak{}latticeScale}}}
\expandafter\newcommand\csname LN@indexEquivY\endcsname{\href{https://github.com/lsernau/grh-lean/blob/72fead102f663d5c7a407964a13dc0cf9d977207/grh-compact/src/New/Trace/GoodIdempotent.lean\#L74-L75}{{\leanfont Data.\allowbreak{}indexEquivY}}}
\expandafter\newcommand\csname LN@existsspaceY\endcsname{\href{https://github.com/lsernau/grh-lean/blob/72fead102f663d5c7a407964a13dc0cf9d977207/grh-compact/src/New/Trace/GoodIdempotent.lean\#L80-L91}{{\leanfont Data.\allowbreak{}exists\_spaceY}}}
\expandafter\newcommand\csname LN@spaceY\endcsname{\href{https://github.com/lsernau/grh-lean/blob/72fead102f663d5c7a407964a13dc0cf9d977207/grh-compact/src/New/Trace/GoodIdempotent.lean\#L95-L95}{{\leanfont Data.\allowbreak{}spaceY}}}
\expandafter\newcommand\csname LN@isStdspaceY\endcsname{\href{https://github.com/lsernau/grh-lean/blob/72fead102f663d5c7a407964a13dc0cf9d977207/grh-compact/src/New/Trace/GoodIdempotent.lean\#L97-L97}{{\leanfont Data.\allowbreak{}isStd\_spaceY}}}
\expandafter\newcommand\csname LN@spaceW\endcsname{\href{https://github.com/lsernau/grh-lean/blob/72fead102f663d5c7a407964a13dc0cf9d977207/grh-compact/src/New/Trace/GoodIdempotent.lean\#L101-L101}{{\leanfont Data.\allowbreak{}spaceW}}}
\expandafter\newcommand\csname LN@principalForms\endcsname{\href{https://github.com/lsernau/grh-lean/blob/72fead102f663d5c7a407964a13dc0cf9d977207/grh-compact/src/New/Trace/GoodIdempotent.lean\#L105-L106}{{\leanfont Data.\allowbreak{}principalForms}}}
\expandafter\newcommand\csname LN@principalRep\endcsname{\href{https://github.com/lsernau/grh-lean/blob/72fead102f663d5c7a407964a13dc0cf9d977207/grh-compact/src/New/Trace/GoodIdempotent.lean\#L111-L111}{{\leanfont Data.\allowbreak{}principalRep}}}
\expandafter\newcommand\csname LN@goodForms\endcsname{\href{https://github.com/lsernau/grh-lean/blob/72fead102f663d5c7a407964a13dc0cf9d977207/grh-compact/src/New/Trace/GoodIdempotent.lean\#L115-L115}{{\leanfont Data.\allowbreak{}goodForms}}}
\expandafter\newcommand\csname LN@rows\endcsname{\href{https://github.com/lsernau/grh-lean/blob/72fead102f663d5c7a407964a13dc0cf9d977207/grh-compact/src/New/Trace/GoodIdempotent.lean\#L118-L118}{{\leanfont Data.\allowbreak{}rows}}}
\expandafter\newcommand\csname LN@rowDim\endcsname{\href{https://github.com/lsernau/grh-lean/blob/72fead102f663d5c7a407964a13dc0cf9d977207/grh-compact/src/New/Trace/GoodIdempotent.lean\#L121-L121}{{\leanfont Data.\allowbreak{}rowDim}}}
\expandafter\newcommand\csname LN@auxPair\endcsname{\href{https://github.com/lsernau/grh-lean/blob/72fead102f663d5c7a407964a13dc0cf9d977207/grh-compact/src/New/Trace/GoodIdempotent.lean\#L125-L130}{{\leanfont Data.\allowbreak{}auxPair}}}
\expandafter\newcommand\csname LN@varsigmaaddwgap\endcsname{\href{https://github.com/lsernau/grh-lean/blob/72fead102f663d5c7a407964a13dc0cf9d977207/grh-compact/src/New/Trace/GoodIdempotent.lean\#L138-L143}{{\leanfont Data.\allowbreak{}varsigma\_add\_w\_gap}}}
\expandafter\newcommand\csname LN@hasGapsignsYparamW\endcsname{\href{https://github.com/lsernau/grh-lean/blob/72fead102f663d5c7a407964a13dc0cf9d977207/grh-compact/src/New/Trace/GoodIdempotent.lean\#L145-L146}{{\leanfont Data.\allowbreak{}hasGap\_signsY\_paramW}}}
\expandafter\newcommand\csname LN@hasGapparamW\endcsname{\href{https://github.com/lsernau/grh-lean/blob/72fead102f663d5c7a407964a13dc0cf9d977207/grh-compact/src/New/Trace/GoodIdempotent.lean\#L148-L149}{{\leanfont Data.\allowbreak{}hasGap\_paramW}}}
\expandafter\newcommand\csname LN@isCuspKernelparamW\endcsname{\href{https://github.com/lsernau/grh-lean/blob/72fead102f663d5c7a407964a13dc0cf9d977207/grh-compact/src/New/Trace/GoodIdempotent.lean\#L154-L155}{{\leanfont Data.\allowbreak{}isCuspKernel\_paramW}}}
\expandafter\newcommand\csname LN@isCentralProje\endcsname{\href{https://github.com/lsernau/grh-lean/blob/72fead102f663d5c7a407964a13dc0cf9d977207/grh-compact/src/New/Trace/GoodIdempotent.lean\#L166-L166}{{\leanfont Data.\allowbreak{}isCentralProj\_e₀}}}
\expandafter\newcommand\csname LN@reflectionSummule\endcsname{\href{https://github.com/lsernau/grh-lean/blob/72fead102f663d5c7a407964a13dc0cf9d977207/grh-compact/src/New/Trace/GoodIdempotent.lean\#L171-L174}{{\leanfont Data.\allowbreak{}reflectionSum\_mul\_e₀}}}
\expandafter\newcommand\csname LN@hasParahoricTypechosenGLRep\endcsname{\href{https://github.com/lsernau/grh-lean/blob/72fead102f663d5c7a407964a13dc0cf9d977207/grh-compact/src/New/Trace/LabelRigidity.lean\#L15-L25}{{\leanfont Data.\allowbreak{}hasParahoricType\_chosenGLRep}}}
\expandafter\newcommand\csname LN@HasRigidLabels\endcsname{\href{https://github.com/lsernau/grh-lean/blob/72fead102f663d5c7a407964a13dc0cf9d977207/grh-compact/src/New/Trace/LabelRigidity.lean\#L29-L30}{{\leanfont Data.\allowbreak{}HasRigidLabels}}}
\expandafter\newcommand\csname LN@hasRigidLabelsspaceW\endcsname{\href{https://github.com/lsernau/grh-lean/blob/72fead102f663d5c7a407964a13dc0cf9d977207/grh-compact/src/New/Trace/LabelRigidity.lean\#L34-L38}{{\leanfont Data.\allowbreak{}hasRigidLabels\_spaceW}}}
\expandafter\newcommand\csname LN@hasRigidLabelsspaceYzero\endcsname{\href{https://github.com/lsernau/grh-lean/blob/72fead102f663d5c7a407964a13dc0cf9d977207/grh-compact/src/New/Trace/LabelRigidity.lean\#L42-L53}{{\leanfont Data.\allowbreak{}hasRigidLabels\_spaceY\_zero}}}
\expandafter\newcommand\csname LN@ezero\endcsname{\href{https://github.com/lsernau/grh-lean/blob/72fead102f663d5c7a407964a13dc0cf9d977207/grh-compact/src/New/Seesaw/IndexTypes.lean\#L314-L314}{{\leanfont Data.\allowbreak{}e₀}}}
\expandafter\newcommand\csname LN@heckeAlgebra\endcsname{\href{https://github.com/lsernau/grh-lean/blob/72fead102f663d5c7a407964a13dc0cf9d977207/grh-compact/src/New/Selmer/FitTwists.lean\#L122-L127}{{\leanfont Data.\allowbreak{}heckeAlgebra}}}
\expandafter\newcommand\csname LN@heckeOperator\endcsname{\href{https://github.com/lsernau/grh-lean/blob/72fead102f663d5c7a407964a13dc0cf9d977207/grh-compact/src/New/Selmer/FitTwists.lean\#L131-L136}{{\leanfont Data.\allowbreak{}heckeOperator}}}
\expandafter\newcommand\csname LN@eigenPart\endcsname{\href{https://github.com/lsernau/grh-lean/blob/72fead102f663d5c7a407964a13dc0cf9d977207/grh-compact/src/New/Selmer/FitTwists.lean\#L140-L140}{{\leanfont Data.\allowbreak{}eigenPart}}}
\expandafter\newcommand\csname LN@occurringLabels\endcsname{\href{https://github.com/lsernau/grh-lean/blob/72fead102f663d5c7a407964a13dc0cf9d977207/grh-compact/src/New/Selmer/FitTwists.lean\#L152-L153}{{\leanfont Data.\allowbreak{}occurringLabels}}}
\expandafter\newcommand\csname LN@heckeOnGood\endcsname{\href{https://github.com/lsernau/grh-lean/blob/72fead102f663d5c7a407964a13dc0cf9d977207/grh-compact/src/New/Selmer/FitTwists.lean\#L177-L179}{{\leanfont Data.\allowbreak{}heckeOnGood}}}
\expandafter\newcommand\csname LN@traceComparison\endcsname{\href{https://github.com/lsernau/grh-lean/blob/72fead102f663d5c7a407964a13dc0cf9d977207/grh-compact/src/Lib/LocalGL/TraceFormula/TraceComparison.lean\#L195-L283}{{\leanfont Space.\allowbreak{}exists\_trace\_eq}}}
\expandafter\newcommand\csname LQ@traceComparison\endcsname{\leanquoteline{0}{theorem exists\_trace\_eq [Nonempty n] [StablyUnram E]}\leanquoteline{4}{l H (j : gapThreshold n (10 * card n \textasciicircum{} 5) ≤ H) : ∃ t, ∀ \{s\} (V : Space E n}\leanquoteline{2}{s) \{ι\} [Fintype ι] (b : OrthonormalBasis ι ℂ (levelSpace (V.formRep (typeKernel l s)) K)),}\leanquoteline{2}{IsCuspKernel B (typeKernel l s) →}\leanquoteline{2}{(∀ w, (toHCParam s (l w)).HasGap H) → ∑ i, ⟪(b i).1, V.formRep \_ h (b i)⟫\_ℂ = t}}
\expandafter\newcommand\csname LN@isCuspKerneltypeKernel\endcsname{\href{https://github.com/lsernau/grh-lean/blob/72fead102f663d5c7a407964a13dc0cf9d977207/grh-compact/src/Lib/Adelic/CuspForms/RealConvolution.lean\#L84-L92}{{\leanfont isCuspKernel\_typeKernel}}}
\expandafter\newcommand\csname LN@finitelevelSpace\endcsname{\href{https://github.com/lsernau/grh-lean/blob/72fead102f663d5c7a407964a13dc0cf9d977207/grh-compact/src/Lib/Adelic/CuspForms/ReproducingKernel.lean\#L54-L64}{{\leanfont Space.\allowbreak{}finite\_levelSpace}}}
\expandafter\newcommand\csname LN@existsframegradedGram\endcsname{\href{https://github.com/lsernau/grh-lean/blob/72fead102f663d5c7a407964a13dc0cf9d977207/grh-compact/src/New/Seesaw/FramesForY.lean\#L29-L36}{{\leanfont exists\_frame\_gradedGram}}}
\expandafter\newcommand\csname LN@isUnimodulargradedLatticeBasisgram\endcsname{\href{https://github.com/lsernau/grh-lean/blob/72fead102f663d5c7a407964a13dc0cf9d977207/grh-compact/src/New/Seesaw/FramesForY.lean\#L15-L24}{{\leanfont isUnimodular\_gradedLatticeBasis\_gram}}}
\expandafter\newcommand\csname LN@existsorthonormal\endcsname{\href{https://github.com/lsernau/grh-lean/blob/72fead102f663d5c7a407964a13dc0cf9d977207/grh-compact/src/Lib/Adelic/Witt/SpaceOperations.lean\#L83-L87}{{\leanfont exists\_orthonormal\_Ô}}}
\expandafter\newcommand\csname LN@isCentralProjsumtupleIdem\endcsname{\href{https://github.com/lsernau/grh-lean/blob/72fead102f663d5c7a407964a13dc0cf9d977207/grh-compact/src/Lib/LocalGL/Green/LabelIdempotents.lean\#L124-L125}{{\leanfont isCentralProj\_sum\_tupleIdem}}}
\expandafter\newcommand\csname LN@reflectionSummullabelIdem\endcsname{\href{https://github.com/lsernau/grh-lean/blob/72fead102f663d5c7a407964a13dc0cf9d977207/grh-compact/src/Lib/LocalGL/Green/LabelIdempotents.lean\#L77-L101}{{\leanfont reflectionSum\_mul\_labelIdem}}}
\expandafter\newcommand\csname LN@piTensorElemapplyeqzero\endcsname{\href{https://github.com/lsernau/grh-lean/blob/72fead102f663d5c7a407964a13dc0cf9d977207/grh-compact/src/Lib/General/Groups/CentralProjections.lean\#L37-L39}{{\leanfont piTensorElem\_apply\_eq\_zero}}}
\expandafter\newcommand\csname LN@charpolyeqmuloftracepowadd\endcsname{\href{https://github.com/lsernau/grh-lean/blob/72fead102f663d5c7a407964a13dc0cf9d977207/grh-compact/src/Lib/GaloisRep/Pseudocharacter/FibreSum.lean\#L70-L76}{{\leanfont charpoly\_eq\_mul\_of\_trace\_pow\_add}}}
\expandafter\newcommand\csname LN@commuteprep\endcsname{\href{https://github.com/lsernau/grh-lean/blob/72fead102f663d5c7a407964a13dc0cf9d977207/grh-compact/src/New/Selmer/FitTwists.lean\#L169-L173}{{\leanfont Data.\allowbreak{}commute\_heckeAlgebra\_principalRep}}}
\expandafter\newcommand\csname LN@iSupeigenPart\endcsname{\href{https://github.com/lsernau/grh-lean/blob/72fead102f663d5c7a407964a13dc0cf9d977207/grh-compact/src/New/Selmer/FitTwists.lean\#L192-L197}{{\leanfont Data.\allowbreak{}iSup\_eigenPart}}}
\expandafter\newcommand\csname LN@finiteeigenPart\endcsname{\href{https://github.com/lsernau/grh-lean/blob/72fead102f663d5c7a407964a13dc0cf9d977207/grh-compact/src/New/Selmer/FitTwists.lean\#L187-L188}{{\leanfont Data.\allowbreak{}finite\_eigenPart}}}
\expandafter\newcommand\csname LN@existsheckeproj\endcsname{\href{https://github.com/lsernau/grh-lean/blob/72fead102f663d5c7a407964a13dc0cf9d977207/grh-compact/src/New/Selmer/FitTwists.lean\#L162-L167}{{\leanfont Data.\allowbreak{}exists\_hecke\_proj\_eigenPart}}}
\expandafter\newcommand\csname LN@occurringLabelssubsingleton\endcsname{\href{https://github.com/lsernau/grh-lean/blob/72fead102f663d5c7a407964a13dc0cf9d977207/grh-compact/src/New/GaloisRep/LabelUniqueness.lean\#L23-L59}{{\leanfont Data.\allowbreak{}occurringLabels\_subsingleton}}}
\expandafter\newcommand\csname LN@auxLift\endcsname{\href{https://github.com/lsernau/grh-lean/blob/72fead102f663d5c7a407964a13dc0cf9d977207/grh-compact/src/New/Selmer/FitTwists.lean\#L226-L227}{{\leanfont Data.\allowbreak{}auxLift}}}
\expandafter\newcommand\csname LN@auxLiftinjective\endcsname{\href{https://github.com/lsernau/grh-lean/blob/72fead102f663d5c7a407964a13dc0cf9d977207/grh-compact/src/New/Exchange/CharacterLiftNorm.lean\#L147-L148}{{\leanfont Data.\allowbreak{}auxLift\_injective}}}
\expandafter\newcommand\csname LN@liftsauxLift\endcsname{\href{https://github.com/lsernau/grh-lean/blob/72fead102f663d5c7a407964a13dc0cf9d977207/grh-compact/src/New/Exchange/CharacterLiftNorm.lean\#L137-L143}{{\leanfont Data.\allowbreak{}lifts\_auxLift}}}
\expandafter\newcommand\csname LN@auxLiftprincipalRep\endcsname{\href{https://github.com/lsernau/grh-lean/blob/72fead102f663d5c7a407964a13dc0cf9d977207/grh-compact/src/New/Exchange/CharacterLiftNorm.lean\#L152-L158}{{\leanfont Data.\allowbreak{}auxLift\_principalRep}}}
\expandafter\newcommand\csname LN@iotaAtLevel\endcsname{\href{https://github.com/lsernau/grh-lean/blob/72fead102f663d5c7a407964a13dc0cf9d977207/grh-compact/src/New/Seesaw/GradingsOfY.lean\#L141-L145}{{\leanfont Data.\allowbreak{}iotaAtLevel}}}
\expandafter\newcommand\csname LN@iotaAtLevelinjective\endcsname{\href{https://github.com/lsernau/grh-lean/blob/72fead102f663d5c7a407964a13dc0cf9d977207/grh-compact/src/New/Exchange/ChainNorms.lean\#L53-L55}{{\leanfont Data.\allowbreak{}iotaAtLevel\_injective}}}
\expandafter\newcommand\csname LN@existsisSphericalWithiota\endcsname{\href{https://github.com/lsernau/grh-lean/blob/72fead102f663d5c7a407964a13dc0cf9d977207/grh-compact/src/New/Exchange/CharacterLiftNorm.lean\#L91-L101}{{\leanfont Data.\allowbreak{}exists\_isSphericalWith\_iota}}}
\expandafter\newcommand\csname LN@algFormRepiota\endcsname{\href{https://github.com/lsernau/grh-lean/blob/72fead102f663d5c7a407964a13dc0cf9d977207/grh-compact/src/New/Seesaw/GradingsOfY.lean\#L111-L131}{{\leanfont Data.\allowbreak{}algFormRep\_iota}}}
\expandafter\newcommand\csname LN@smulchosenGLRepeq\endcsname{\href{https://github.com/lsernau/grh-lean/blob/72fead102f663d5c7a407964a13dc0cf9d977207/grh-compact/src/New/SiegelWeil/AuxPairVanishing.lean\#L37-L49}{{\leanfont AuxPair.\allowbreak{}smul\_chosenGLRep\_eq}}}
\expandafter\newcommand\csname LN@hasGapgridParam\endcsname{\href{https://github.com/lsernau/grh-lean/blob/72fead102f663d5c7a407964a13dc0cf9d977207/grh-compact/src/New/Exchange/GridParameters.lean\#L121-L124}{{\leanfont Data.\allowbreak{}hasGap\_gridParam}}}
\expandafter\newcommand\csname LN@cardneggridSigns\endcsname{\href{https://github.com/lsernau/grh-lean/blob/72fead102f663d5c7a407964a13dc0cf9d977207/grh-compact/src/New/Exchange/GridParameters.lean\#L81-L83}{{\leanfont Data.\allowbreak{}card\_neg\_gridSigns}}}
\expandafter\newcommand\csname LN@omegazpowneone\endcsname{\href{https://github.com/lsernau/grh-lean/blob/72fead102f663d5c7a407964a13dc0cf9d977207/grh-compact/src/New/Data/ChosenPlaceLevels.lean\#L86-L90}{{\leanfont Data.\allowbreak{}omega\_zpow\_ne\_one}}}
\newcommand{\ConcordanceRows}{\crow{existscommonheckerallis}{\href{https://github.com/lsernau/grh-lean/blob/72fead102f663d5c7a407964a13dc0cf9d977207/grh-compact/src/New/Trace/GoodFormTransfer.lean\#L27-L47}{{\leanfont Data.\allowbreak{}exists\_common\_hecke\_rallis}}}{GoodFormTransfer.lean}{27--47}
\crow{traceheckeOnGoodeq}{\href{https://github.com/lsernau/grh-lean/blob/72fead102f663d5c7a407964a13dc0cf9d977207/grh-compact/src/New/Trace/GoodFormTransfer.lean\#L56-L89}{{\leanfont Data.\allowbreak{}trace\_heckeOnGood\_eq}}}{GoodFormTransfer.lean}{56--89}
\crow{detheckeOnGoodeq}{\href{https://github.com/lsernau/grh-lean/blob/72fead102f663d5c7a407964a13dc0cf9d977207/grh-compact/src/New/Trace/GoodFormTransfer.lean\#L94-L98}{{\leanfont Data.\allowbreak{}det\_heckeOnGood\_eq}}}{GoodFormTransfer.lean}{94--98}
\crow{goodFormsinfeigenParteqbot}{\href{https://github.com/lsernau/grh-lean/blob/72fead102f663d5c7a407964a13dc0cf9d977207/grh-compact/src/New/Trace/GoodFormTransfer.lean\#L103-L110}{{\leanfont Data.\allowbreak{}goodForms\_inf\_eigenPart\_eq\_bot}}}{GoodFormTransfer.lean}{103--110}
\crow{finrankrowseqfinrankgoodForms}{\href{https://github.com/lsernau/grh-lean/blob/72fead102f663d5c7a407964a13dc0cf9d977207/grh-compact/src/New/Trace/GoodFormTransfer.lean\#L116-L118}{{\leanfont Data.\allowbreak{}finrank\_rows\_eq\_finrank\_goodForms}}}{GoodFormTransfer.lean}{116--118}
\crow{besselIndexEquiv}{\href{https://github.com/lsernau/grh-lean/blob/72fead102f663d5c7a407964a13dc0cf9d977207/grh-compact/src/New/Trace/GoodIdempotent.lean\#L16-L19}{{\leanfont Shape.\allowbreak{}besselIndexEquiv}}}{GoodIdempotent.lean}{16--19}
\crow{meqtwomulvarrhoaddtwo}{\href{https://github.com/lsernau/grh-lean/blob/72fead102f663d5c7a407964a13dc0cf9d977207/grh-compact/src/New/Trace/GoodIdempotent.lean\#L27-L27}{{\leanfont Data.\allowbreak{}m\_eq\_two\_mul\_varrho\_add\_two}}}{GoodIdempotent.lean}{27--27}
\crow{levelNormmemandnotMemofmemQ}{\href{https://github.com/lsernau/grh-lean/blob/72fead102f663d5c7a407964a13dc0cf9d977207/grh-compact/src/New/Trace/GoodIdempotent.lean\#L34-L36}{{\leanfont Data.\allowbreak{}levelNorm\_mem\_and\_notMem\_of\_mem\_Q}}}{GoodIdempotent.lean}{34--36}
\crow{levelPlaces}{\href{https://github.com/lsernau/grh-lean/blob/72fead102f663d5c7a407964a13dc0cf9d977207/grh-compact/src/New/Trace/GoodIdempotent.lean\#L40-L42}{{\leanfont Data.\allowbreak{}levelPlaces}}}{GoodIdempotent.lean}{40--42}
\crow{decl}{\href{https://github.com/lsernau/grh-lean/blob/72fead102f663d5c7a407964a13dc0cf9d977207/grh-compact/src/New/Trace/GoodIdempotent.lean\#L46-L46}{{\leanfont Data.\allowbreak{}𝔪}}}{GoodIdempotent.lean}{46--46}
\crow{nonemptymodulusplaces}{\href{https://github.com/lsernau/grh-lean/blob/72fead102f663d5c7a407964a13dc0cf9d977207/grh-compact/src/New/Trace/GoodIdempotent.lean\#L48-L50}{{\leanfont Data.\allowbreak{}nonempty\_modulus\_places}}}{GoodIdempotent.lean}{48--50}
\crow{chosenIdele}{\href{https://github.com/lsernau/grh-lean/blob/72fead102f663d5c7a407964a13dc0cf9d977207/grh-compact/src/New/Trace/GoodIdempotent.lean\#L54-L55}{{\leanfont Data.\allowbreak{}chosenIdele}}}{GoodIdempotent.lean}{54--55}
\crow{ideleA}{\href{https://github.com/lsernau/grh-lean/blob/72fead102f663d5c7a407964a13dc0cf9d977207/grh-compact/src/New/Trace/GoodIdempotent.lean\#L59-L59}{{\leanfont Data.\allowbreak{}ideleA}}}{GoodIdempotent.lean}{59--59}
\crow{piIdele}{\href{https://github.com/lsernau/grh-lean/blob/72fead102f663d5c7a407964a13dc0cf9d977207/grh-compact/src/New/Trace/GoodIdempotent.lean\#L62-L62}{{\leanfont Data.\allowbreak{}piIdele}}}{GoodIdempotent.lean}{62--62}
\crow{piIdelespec}{\href{https://github.com/lsernau/grh-lean/blob/72fead102f663d5c7a407964a13dc0cf9d977207/grh-compact/src/New/Trace/GoodIdempotent.lean\#L64-L66}{{\leanfont Data.\allowbreak{}piIdele\_spec}}}{GoodIdempotent.lean}{64--66}
\crow{latticeScale}{\href{https://github.com/lsernau/grh-lean/blob/72fead102f663d5c7a407964a13dc0cf9d977207/grh-compact/src/New/Trace/GoodIdempotent.lean\#L70-L70}{{\leanfont Data.\allowbreak{}latticeScale}}}{GoodIdempotent.lean}{70--70}
\crow{indexEquivY}{\href{https://github.com/lsernau/grh-lean/blob/72fead102f663d5c7a407964a13dc0cf9d977207/grh-compact/src/New/Trace/GoodIdempotent.lean\#L74-L75}{{\leanfont Data.\allowbreak{}indexEquivY}}}{GoodIdempotent.lean}{74--75}
\crow{existsspaceY}{\href{https://github.com/lsernau/grh-lean/blob/72fead102f663d5c7a407964a13dc0cf9d977207/grh-compact/src/New/Trace/GoodIdempotent.lean\#L80-L91}{{\leanfont Data.\allowbreak{}exists\_spaceY}}}{GoodIdempotent.lean}{80--91}
\crow{spaceY}{\href{https://github.com/lsernau/grh-lean/blob/72fead102f663d5c7a407964a13dc0cf9d977207/grh-compact/src/New/Trace/GoodIdempotent.lean\#L95-L95}{{\leanfont Data.\allowbreak{}spaceY}}}{GoodIdempotent.lean}{95--95}
\crow{isStdspaceY}{\href{https://github.com/lsernau/grh-lean/blob/72fead102f663d5c7a407964a13dc0cf9d977207/grh-compact/src/New/Trace/GoodIdempotent.lean\#L97-L97}{{\leanfont Data.\allowbreak{}isStd\_spaceY}}}{GoodIdempotent.lean}{97--97}
\crow{spaceW}{\href{https://github.com/lsernau/grh-lean/blob/72fead102f663d5c7a407964a13dc0cf9d977207/grh-compact/src/New/Trace/GoodIdempotent.lean\#L101-L101}{{\leanfont Data.\allowbreak{}spaceW}}}{GoodIdempotent.lean}{101--101}
\crow{principalForms}{\href{https://github.com/lsernau/grh-lean/blob/72fead102f663d5c7a407964a13dc0cf9d977207/grh-compact/src/New/Trace/GoodIdempotent.lean\#L105-L106}{{\leanfont Data.\allowbreak{}principalForms}}}{GoodIdempotent.lean}{105--106}
\crow{principalRep}{\href{https://github.com/lsernau/grh-lean/blob/72fead102f663d5c7a407964a13dc0cf9d977207/grh-compact/src/New/Trace/GoodIdempotent.lean\#L111-L111}{{\leanfont Data.\allowbreak{}principalRep}}}{GoodIdempotent.lean}{111--111}
\crow{goodForms}{\href{https://github.com/lsernau/grh-lean/blob/72fead102f663d5c7a407964a13dc0cf9d977207/grh-compact/src/New/Trace/GoodIdempotent.lean\#L115-L115}{{\leanfont Data.\allowbreak{}goodForms}}}{GoodIdempotent.lean}{115--115}
\crow{rows}{\href{https://github.com/lsernau/grh-lean/blob/72fead102f663d5c7a407964a13dc0cf9d977207/grh-compact/src/New/Trace/GoodIdempotent.lean\#L118-L118}{{\leanfont Data.\allowbreak{}rows}}}{GoodIdempotent.lean}{118--118}
\crow{rowDim}{\href{https://github.com/lsernau/grh-lean/blob/72fead102f663d5c7a407964a13dc0cf9d977207/grh-compact/src/New/Trace/GoodIdempotent.lean\#L121-L121}{{\leanfont Data.\allowbreak{}rowDim}}}{GoodIdempotent.lean}{121--121}
\crow{auxPair}{\href{https://github.com/lsernau/grh-lean/blob/72fead102f663d5c7a407964a13dc0cf9d977207/grh-compact/src/New/Trace/GoodIdempotent.lean\#L125-L130}{{\leanfont Data.\allowbreak{}auxPair}}}{GoodIdempotent.lean}{125--130}
\crow{varsigmaaddwgap}{\href{https://github.com/lsernau/grh-lean/blob/72fead102f663d5c7a407964a13dc0cf9d977207/grh-compact/src/New/Trace/GoodIdempotent.lean\#L138-L143}{{\leanfont Data.\allowbreak{}varsigma\_add\_w\_gap}}}{GoodIdempotent.lean}{138--143}
\crow{hasGapsignsYparamW}{\href{https://github.com/lsernau/grh-lean/blob/72fead102f663d5c7a407964a13dc0cf9d977207/grh-compact/src/New/Trace/GoodIdempotent.lean\#L145-L146}{{\leanfont Data.\allowbreak{}hasGap\_signsY\_paramW}}}{GoodIdempotent.lean}{145--146}
\crow{hasGapparamW}{\href{https://github.com/lsernau/grh-lean/blob/72fead102f663d5c7a407964a13dc0cf9d977207/grh-compact/src/New/Trace/GoodIdempotent.lean\#L148-L149}{{\leanfont Data.\allowbreak{}hasGap\_paramW}}}{GoodIdempotent.lean}{148--149}
\crow{isCuspKernelparamW}{\href{https://github.com/lsernau/grh-lean/blob/72fead102f663d5c7a407964a13dc0cf9d977207/grh-compact/src/New/Trace/GoodIdempotent.lean\#L154-L155}{{\leanfont Data.\allowbreak{}isCuspKernel\_paramW}}}{GoodIdempotent.lean}{154--155}
\crow{isCentralProje}{\href{https://github.com/lsernau/grh-lean/blob/72fead102f663d5c7a407964a13dc0cf9d977207/grh-compact/src/New/Trace/GoodIdempotent.lean\#L166-L166}{{\leanfont Data.\allowbreak{}isCentralProj\_e₀}}}{GoodIdempotent.lean}{166--166}
\crow{reflectionSummule}{\href{https://github.com/lsernau/grh-lean/blob/72fead102f663d5c7a407964a13dc0cf9d977207/grh-compact/src/New/Trace/GoodIdempotent.lean\#L171-L174}{{\leanfont Data.\allowbreak{}reflectionSum\_mul\_e₀}}}{GoodIdempotent.lean}{171--174}
\crow{hasParahoricTypechosenGLRep}{\href{https://github.com/lsernau/grh-lean/blob/72fead102f663d5c7a407964a13dc0cf9d977207/grh-compact/src/New/Trace/LabelRigidity.lean\#L15-L25}{{\leanfont Data.\allowbreak{}hasParahoricType\_chosenGLRep}}}{LabelRigidity.lean}{15--25}
\crow{HasRigidLabels}{\href{https://github.com/lsernau/grh-lean/blob/72fead102f663d5c7a407964a13dc0cf9d977207/grh-compact/src/New/Trace/LabelRigidity.lean\#L29-L30}{{\leanfont Data.\allowbreak{}HasRigidLabels}}}{LabelRigidity.lean}{29--30}
\crow{hasRigidLabelsspaceW}{\href{https://github.com/lsernau/grh-lean/blob/72fead102f663d5c7a407964a13dc0cf9d977207/grh-compact/src/New/Trace/LabelRigidity.lean\#L34-L38}{{\leanfont Data.\allowbreak{}hasRigidLabels\_spaceW}}}{LabelRigidity.lean}{34--38}
\crow{hasRigidLabelsspaceYzero}{\href{https://github.com/lsernau/grh-lean/blob/72fead102f663d5c7a407964a13dc0cf9d977207/grh-compact/src/New/Trace/LabelRigidity.lean\#L42-L53}{{\leanfont Data.\allowbreak{}hasRigidLabels\_spaceY\_zero}}}{LabelRigidity.lean}{42--53}
\crow{ezero}{\href{https://github.com/lsernau/grh-lean/blob/72fead102f663d5c7a407964a13dc0cf9d977207/grh-compact/src/New/Seesaw/IndexTypes.lean\#L314-L314}{{\leanfont Data.\allowbreak{}e₀}}}{IndexTypes.lean}{314--314}
\crow{heckeAlgebra}{\href{https://github.com/lsernau/grh-lean/blob/72fead102f663d5c7a407964a13dc0cf9d977207/grh-compact/src/New/Selmer/FitTwists.lean\#L122-L127}{{\leanfont Data.\allowbreak{}heckeAlgebra}}}{FitTwists.lean}{122--127}
\crow{heckeOperator}{\href{https://github.com/lsernau/grh-lean/blob/72fead102f663d5c7a407964a13dc0cf9d977207/grh-compact/src/New/Selmer/FitTwists.lean\#L131-L136}{{\leanfont Data.\allowbreak{}heckeOperator}}}{FitTwists.lean}{131--136}
\crow{eigenPart}{\href{https://github.com/lsernau/grh-lean/blob/72fead102f663d5c7a407964a13dc0cf9d977207/grh-compact/src/New/Selmer/FitTwists.lean\#L140-L140}{{\leanfont Data.\allowbreak{}eigenPart}}}{FitTwists.lean}{140--140}
\crow{occurringLabels}{\href{https://github.com/lsernau/grh-lean/blob/72fead102f663d5c7a407964a13dc0cf9d977207/grh-compact/src/New/Selmer/FitTwists.lean\#L152-L153}{{\leanfont Data.\allowbreak{}occurringLabels}}}{FitTwists.lean}{152--153}
\crow{heckeOnGood}{\href{https://github.com/lsernau/grh-lean/blob/72fead102f663d5c7a407964a13dc0cf9d977207/grh-compact/src/New/Selmer/FitTwists.lean\#L177-L179}{{\leanfont Data.\allowbreak{}heckeOnGood}}}{FitTwists.lean}{177--179}
\crow{traceComparison}{\href{https://github.com/lsernau/grh-lean/blob/72fead102f663d5c7a407964a13dc0cf9d977207/grh-compact/src/Lib/LocalGL/TraceFormula/TraceComparison.lean\#L195-L283}{{\leanfont Space.\allowbreak{}exists\_trace\_eq}}}{TraceComparison.lean}{195--283}
\crow{isCuspKerneltypeKernel}{\href{https://github.com/lsernau/grh-lean/blob/72fead102f663d5c7a407964a13dc0cf9d977207/grh-compact/src/Lib/Adelic/CuspForms/RealConvolution.lean\#L84-L92}{{\leanfont isCuspKernel\_typeKernel}}}{RealConvolution.lean}{84--92}
\crow{finitelevelSpace}{\href{https://github.com/lsernau/grh-lean/blob/72fead102f663d5c7a407964a13dc0cf9d977207/grh-compact/src/Lib/Adelic/CuspForms/ReproducingKernel.lean\#L54-L64}{{\leanfont Space.\allowbreak{}finite\_levelSpace}}}{ReproducingKernel.lean}{54--64}
\crow{existsframegradedGram}{\href{https://github.com/lsernau/grh-lean/blob/72fead102f663d5c7a407964a13dc0cf9d977207/grh-compact/src/New/Seesaw/FramesForY.lean\#L29-L36}{{\leanfont exists\_frame\_gradedGram}}}{FramesForY.lean}{29--36}
\crow{isUnimodulargradedLatticeBasisgram}{\href{https://github.com/lsernau/grh-lean/blob/72fead102f663d5c7a407964a13dc0cf9d977207/grh-compact/src/New/Seesaw/FramesForY.lean\#L15-L24}{{\leanfont isUnimodular\_gradedLatticeBasis\_gram}}}{FramesForY.lean}{15--24}
\crow{existsorthonormal}{\href{https://github.com/lsernau/grh-lean/blob/72fead102f663d5c7a407964a13dc0cf9d977207/grh-compact/src/Lib/Adelic/Witt/SpaceOperations.lean\#L83-L87}{{\leanfont exists\_orthonormal\_Ô}}}{SpaceOperations.lean}{83--87}
\crow{isCentralProjsumtupleIdem}{\href{https://github.com/lsernau/grh-lean/blob/72fead102f663d5c7a407964a13dc0cf9d977207/grh-compact/src/Lib/LocalGL/Green/LabelIdempotents.lean\#L124-L125}{{\leanfont isCentralProj\_sum\_tupleIdem}}}{LabelIdempotents.lean}{124--125}
\crow{reflectionSummullabelIdem}{\href{https://github.com/lsernau/grh-lean/blob/72fead102f663d5c7a407964a13dc0cf9d977207/grh-compact/src/Lib/LocalGL/Green/LabelIdempotents.lean\#L77-L101}{{\leanfont reflectionSum\_mul\_labelIdem}}}{LabelIdempotents.lean}{77--101}
\crow{piTensorElemapplyeqzero}{\href{https://github.com/lsernau/grh-lean/blob/72fead102f663d5c7a407964a13dc0cf9d977207/grh-compact/src/Lib/General/Groups/CentralProjections.lean\#L37-L39}{{\leanfont piTensorElem\_apply\_eq\_zero}}}{CentralProjections.lean}{37--39}
\crow{charpolyeqmuloftracepowadd}{\href{https://github.com/lsernau/grh-lean/blob/72fead102f663d5c7a407964a13dc0cf9d977207/grh-compact/src/Lib/GaloisRep/Pseudocharacter/FibreSum.lean\#L70-L76}{{\leanfont charpoly\_eq\_mul\_of\_trace\_pow\_add}}}{FibreSum.lean}{70--76}
\crow{commuteprep}{\href{https://github.com/lsernau/grh-lean/blob/72fead102f663d5c7a407964a13dc0cf9d977207/grh-compact/src/New/Selmer/FitTwists.lean\#L169-L173}{{\leanfont Data.\allowbreak{}commute\_heckeAlgebra\_principalRep}}}{FitTwists.lean}{169--173}
\crow{iSupeigenPart}{\href{https://github.com/lsernau/grh-lean/blob/72fead102f663d5c7a407964a13dc0cf9d977207/grh-compact/src/New/Selmer/FitTwists.lean\#L192-L197}{{\leanfont Data.\allowbreak{}iSup\_eigenPart}}}{FitTwists.lean}{192--197}
\crow{finiteeigenPart}{\href{https://github.com/lsernau/grh-lean/blob/72fead102f663d5c7a407964a13dc0cf9d977207/grh-compact/src/New/Selmer/FitTwists.lean\#L187-L188}{{\leanfont Data.\allowbreak{}finite\_eigenPart}}}{FitTwists.lean}{187--188}
\crow{existsheckeproj}{\href{https://github.com/lsernau/grh-lean/blob/72fead102f663d5c7a407964a13dc0cf9d977207/grh-compact/src/New/Selmer/FitTwists.lean\#L162-L167}{{\leanfont Data.\allowbreak{}exists\_hecke\_proj\_eigenPart}}}{FitTwists.lean}{162--167}
\crow{occurringLabelssubsingleton}{\href{https://github.com/lsernau/grh-lean/blob/72fead102f663d5c7a407964a13dc0cf9d977207/grh-compact/src/New/GaloisRep/LabelUniqueness.lean\#L23-L59}{{\leanfont Data.\allowbreak{}occurringLabels\_subsingleton}}}{LabelUniqueness.lean}{23--59}
\crow{auxLift}{\href{https://github.com/lsernau/grh-lean/blob/72fead102f663d5c7a407964a13dc0cf9d977207/grh-compact/src/New/Selmer/FitTwists.lean\#L226-L227}{{\leanfont Data.\allowbreak{}auxLift}}}{FitTwists.lean}{226--227}
\crow{auxLiftinjective}{\href{https://github.com/lsernau/grh-lean/blob/72fead102f663d5c7a407964a13dc0cf9d977207/grh-compact/src/New/Exchange/CharacterLiftNorm.lean\#L147-L148}{{\leanfont Data.\allowbreak{}auxLift\_injective}}}{CharacterLiftNorm.lean}{147--148}
\crow{liftsauxLift}{\href{https://github.com/lsernau/grh-lean/blob/72fead102f663d5c7a407964a13dc0cf9d977207/grh-compact/src/New/Exchange/CharacterLiftNorm.lean\#L137-L143}{{\leanfont Data.\allowbreak{}lifts\_auxLift}}}{CharacterLiftNorm.lean}{137--143}
\crow{auxLiftprincipalRep}{\href{https://github.com/lsernau/grh-lean/blob/72fead102f663d5c7a407964a13dc0cf9d977207/grh-compact/src/New/Exchange/CharacterLiftNorm.lean\#L152-L158}{{\leanfont Data.\allowbreak{}auxLift\_principalRep}}}{CharacterLiftNorm.lean}{152--158}
\crow{iotaAtLevel}{\href{https://github.com/lsernau/grh-lean/blob/72fead102f663d5c7a407964a13dc0cf9d977207/grh-compact/src/New/Seesaw/GradingsOfY.lean\#L141-L145}{{\leanfont Data.\allowbreak{}iotaAtLevel}}}{GradingsOfY.lean}{141--145}
\crow{iotaAtLevelinjective}{\href{https://github.com/lsernau/grh-lean/blob/72fead102f663d5c7a407964a13dc0cf9d977207/grh-compact/src/New/Exchange/ChainNorms.lean\#L53-L55}{{\leanfont Data.\allowbreak{}iotaAtLevel\_injective}}}{ChainNorms.lean}{53--55}
\crow{existsisSphericalWithiota}{\href{https://github.com/lsernau/grh-lean/blob/72fead102f663d5c7a407964a13dc0cf9d977207/grh-compact/src/New/Exchange/CharacterLiftNorm.lean\#L91-L101}{{\leanfont Data.\allowbreak{}exists\_isSphericalWith\_iota}}}{CharacterLiftNorm.lean}{91--101}
\crow{algFormRepiota}{\href{https://github.com/lsernau/grh-lean/blob/72fead102f663d5c7a407964a13dc0cf9d977207/grh-compact/src/New/Seesaw/GradingsOfY.lean\#L111-L131}{{\leanfont Data.\allowbreak{}algFormRep\_iota}}}{GradingsOfY.lean}{111--131}
\crow{smulchosenGLRepeq}{\href{https://github.com/lsernau/grh-lean/blob/72fead102f663d5c7a407964a13dc0cf9d977207/grh-compact/src/New/SiegelWeil/AuxPairVanishing.lean\#L37-L49}{{\leanfont AuxPair.\allowbreak{}smul\_chosenGLRep\_eq}}}{AuxPairVanishing.lean}{37--49}
\crow{hasGapgridParam}{\href{https://github.com/lsernau/grh-lean/blob/72fead102f663d5c7a407964a13dc0cf9d977207/grh-compact/src/New/Exchange/GridParameters.lean\#L121-L124}{{\leanfont Data.\allowbreak{}hasGap\_gridParam}}}{GridParameters.lean}{121--124}
\crow{cardneggridSigns}{\href{https://github.com/lsernau/grh-lean/blob/72fead102f663d5c7a407964a13dc0cf9d977207/grh-compact/src/New/Exchange/GridParameters.lean\#L81-L83}{{\leanfont Data.\allowbreak{}card\_neg\_gridSigns}}}{GridParameters.lean}{81--83}
\crow{omegazpowneone}{\href{https://github.com/lsernau/grh-lean/blob/72fead102f663d5c7a407964a13dc0cf9d977207/grh-compact/src/New/Data/ChosenPlaceLevels.lean\#L86-L90}{{\leanfont Data.\allowbreak{}omega\_zpow\_ne\_one}}}{ChosenPlaceLevels.lean}{86--90}}
% the notations of the defined functions: one slot for each explicit argument of the Lean declaration
\newcommand{\besselIx}[1]{\mathrm{bes}_{#1}}
\newcommand{\levelPlaces}[2]{\mathrm{Lv}(#1,#2)}
\newcommand{\levelMod}[4]{\mathfrak{m}_{#1,#2}^{#3,#4}}
\newcommand{\chosenIdele}[1]{\hat{\pi}_g(#1)}
\newcommand{\ideleA}[1]{\alpha_g(#1)}
\newcommand{\piIdele}[1]{\hat{\pi}_N(#1)}
\newcommand{\latticeScale}[2]{\mathrm{sc}_{#2}(#1)}
\newcommand{\indexY}[3]{\mathrm{ix}_{#1,#2}^{\,#3}}
\newcommand{\spaceY}[2]{Y_{#2}(#1)}
\newcommand{\spaceW}[2]{W_{#2}(#1)}
\newcommand{\Prin}[2]{\mathrm{Prin}(#1,#2)}
\newcommand{\prep}[2]{\mathrm{rep}_{#1,#2}}
\newcommand{\Good}[2]{\mathrm{Good}(#1,#2)}
\newcommand{\Row}[1]{\mathrm{Row}(#1)}
\newcommand{\rowDim}[1]{\dim\mathrm{Row}(#1)}
\newcommand{\auxPair}[3]{\mathrm{Aux}_{#3}(#1,#2)}
\newcommand{\Rigid}[2]{labels are rigid for~$#1$ on~$#2$}
\newcommand{\UnitMembers}{46}
\newcommand{\UnitInNote}{37}
\newcommand{\UnitProofs}{16}
\newcommand{\UnitSteps}{99}
\newcommand{\UnitAllSteps}{126}


\linespread{1.08}                       % 11 pt on 14.7 pt
\setlength{\parindent}{0pt}
\setlength{\parskip}{5.5pt plus 1.5pt minus 0.5pt}
\setlength{\abovedisplayskip}{7.5pt plus 2pt minus 3pt}
\setlength{\belowdisplayskip}{7.5pt plus 2pt minus 3pt}
\setlength{\abovedisplayshortskip}{3pt plus 2pt}
\setlength{\belowdisplayshortskip}{6pt plus 2pt minus 2pt}
\emergencystretch=1.5em
\clubpenalty=10000 \widowpenalty=10000 \displaywidowpenalty=10000
\relpenalty=7000 \binoppenalty=9000
\raggedbottom
\pagestyle{fancy}\fancyhf{}\renewcommand{\headrulewidth}{0pt}
\fancyfoot[C]{\sffamily\fontsize{7.6}{9}\selectfont\color{inkfaint}\thepage\,/\,\pageref*{LastPage}}
\setlist{topsep=2pt,partopsep=0pt,parsep=0pt,itemsep=4pt,leftmargin=2.1em,labelsep=0.55em}

\newcommand{\Q}{\mathbb{Q}}
\newcommand{\Z}{\mathbb{Z}}
\newcommand{\F}{\mathbb{F}}
\newcommand{\len}{\operatorname{len}}
\newcommand{\cok}{\operatorname{cok}}
\newcommand{\End}{\operatorname{End}}
\newcommand{\Hom}{\operatorname{Hom}}
\newcommand{\Gal}{\operatorname{Gal}}
\newcommand{\diag}{\operatorname{diag}}
\newcommand{\kw}[1]{\emph{#1}}
\newcommand{\m}[1]{\mbox{$#1$}}        % a short formula that must not be broken across lines
\newcommand{\head}[1]{\par\addvspace{17pt plus 4pt minus 2pt}\Needspace{6\baselineskip}%
  \noindent{\fontsize{13.2}{16}\selectfont\bfseries\color{ink}#1}\par\nopagebreak\vspace{-1pt}}
\newcommand{\sect}[2]{\refstepcounter{section}\label{#2}\par\addvspace{19pt plus 4pt minus 2pt}\Needspace{7\baselineskip}%
  \noindent\llap{{\fontsize{13.2}{16}\selectfont\bfseries\color{inkfaint}\thesection}\hspace{9pt}}%
  {\fontsize{13.2}{16}\selectfont\bfseries\color{ink}#1}\par\nopagebreak\vspace{-1pt}}
\newcommand{\subhead}[1]{\par\addvspace{9pt plus 2pt}\Needspace{4\baselineskip}\noindent\textbf{#1}\enspace\ignorespaces}
% where a statement lives in the Lean archive
\newcommand{\leantag}{{\sffamily\bfseries\fontsize{6.3}{8}\selectfont\color{leancol}\addfontfeatures{LetterSpace=7}LEAN}}
\newcommand{\sep}{\unskip\nobreak\hspace{5pt}{\color{inkfaint}\textperiodcentered}\hspace{5pt plus 2pt}\ignorespaces}
\newcommand{\leanline}[1]{\par\nopagebreak\vspace{-2.5pt}{\fontsize{8.9}{12.4}\selectfont\color{inksoft}\raggedright
  \noindent\llap{\leantag\hspace{9pt}}#1\par}\vspace{1pt}}
% a displayed statement in a tinted box: #1 colour, #2 text
\newcommand{\statement}[2]{\par\addvspace{7pt}\begingroup\setlength{\fboxsep}{8.5pt}\setlength{\fboxrule}{0.6pt}%
  \noindent\fcolorbox{#1!30!white}{#1!3!white}{\parbox{\dimexpr\linewidth-2\fboxsep-2\fboxrule\relax}{#2}}%
  \endgroup\par\addvspace{7pt}}
\newcommand{\nt}[2]{\par\noindent\hangindent=1.1em\hangafter=1 #1\enspace{\color{inksoft}#2}\par\vspace{1pt}}
\newcommand{\legdot}[1]{\tikz[baseline=-0.62ex]\fill[#1] (0,0) circle (2.5pt);\hspace{3pt}}
\newcommand{\crow}[4]{\csname CS@#1\endcsname & #2 & \textcolor{inksoft}{#3} & \textcolor{inksoft}{#4}\\}
\newcommand{\leanquoteline}[2]{\par\noindent\hangindent=\dimexpr#1\fontdimen2\font+2.2em\relax\hangafter=1
  \hspace*{#1\fontdimen2\font}#2}
\newcommand{\quoted}[1]{\par\vspace{3pt}{\ttfamily\fontsize{7.5}{10.4}\selectfont\color{leancol}\raggedright
  \setlength{\parskip}{0pt}\leanquote{#1}\par}\vspace{4pt}}
\newcommand{\bx}[3]{{\sffamily\bfseries\fontsize{7.3}{9.2}\selectfont\color{#1!82!black}#2}\\[1.5pt]#3}
\newcommand{\cs}[2]{\expandafter\newcommand\csname CS@#1\endcsname{#2}}

% theorems
\newtheoremstyle{plainx}{7pt plus 2pt}{7pt plus 2pt}{\itshape}{}{\bfseries}{.}{0.5em}{}
\newtheoremstyle{uprightx}{7pt plus 2pt}{7pt plus 2pt}{}{}{\bfseries}{.}{0.5em}{}
\theoremstyle{plainx}
\newtheorem{theorem}{Theorem}[section]
\newtheorem{lemma}[theorem]{Lemma}
\newtheorem{proposition}[theorem]{Proposition}
\newtheorem{corollary}[theorem]{Corollary}
\theoremstyle{uprightx}
\newtheorem{definition}[theorem]{Definition}
\newtheorem{setting}[theorem]{Setting}          % hypotheses shared by several statements; each of them cites it by number
\newtheorem{remark}[theorem]{Remark}
\newtheorem{example}[theorem]{Example}
\renewcommand{\qedsymbol}{{\setlength{\fboxsep}{0pt}\setlength{\fboxrule}{0.5pt}\fbox{\rule{0pt}{0.46em}\hspace{0.46em}}}}
% what holds throughout the note: the ambient objects only (implicit and instance arguments), never a hypothesis
\newenvironment{conventions}{\par\addvspace{5pt}\noindent\textbf{Conventions.}\enspace\ignorespaces}{\par\addvspace{3pt}}
\makeatletter
\renewenvironment{proof}[1][Proof]{\par\pushQED{\qed}\normalfont\topsep0pt\partopsep0pt
  \trivlist\item[\hskip\labelsep\itshape #1.]\ignorespaces}{\popQED\endtrivlist\@endpefalse}
\makeatother

% where each declaration of the concordance is treated: \cs{key}{place}, the place naming a \label or a lettered item
\cs{traceheckeOnGoodeq}{Theorem 1}
\cs{detheckeOnGoodeq}{Theorem 2}
\cs{goodFormsinfeigenParteqbot}{Theorem 2}
\cs{finrankrowseqfinrankgoodForms}{Theorem 2}
\cs{hasRigidLabelsspaceW}{Theorem 3}
\cs{hasRigidLabelsspaceYzero}{Theorem 3}
\cs{besselIndexEquiv}{Definition~\ref{d:index}}
\cs{indexEquivY}{Definition~\ref{d:index}}
\cs{meqtwomulvarrhoaddtwo}{Lemma~\ref{l:m}}
\cs{chosenIdele}{Definition~\ref{d:ideles}}
\cs{ideleA}{Definition~\ref{d:ideles}}
\cs{piIdele}{Definition~\ref{d:ideles}}
\cs{latticeScale}{Definition~\ref{d:ideles}}
\cs{piIdelespec}{Lemma~\ref{l:pi}}
\cs{existsspaceY}{Lemma~\ref{l:Y}}
\cs{spaceY}{Definition~\ref{d:YW}}
\cs{spaceW}{Definition~\ref{d:YW}}
\cs{isStdspaceY}{Lemma~\ref{l:std}}
\cs{principalForms}{Definition~\ref{d:prin}}
\cs{principalRep}{Definition~\ref{d:prin}}
\cs{goodForms}{Definition~\ref{d:good}}
\cs{rows}{Definition~\ref{d:good}}
\cs{rowDim}{Definition~\ref{d:good}}
\cs{HasRigidLabels}{Definition~\ref{d:rigid}}
\cs{varsigmaaddwgap}{Lemma~\ref{l:gap}}
\cs{hasGapsignsYparamW}{Lemma~\ref{l:hasgap}}
\cs{hasGapparamW}{Lemma~\ref{l:hasgap}}
\cs{isCuspKernelparamW}{Lemma~\ref{l:cusp}}
\cs{levelNormmemandnotMemofmemQ}{Lemma~\ref{l:Q}}
\cs{levelPlaces}{Definition~\ref{d:levels}}
\cs{decl}{Definition~\ref{d:levels}}
\cs{nonemptymodulusplaces}{Lemma~\ref{l:nonempty}}
\cs{auxPair}{Definition~\ref{d:aux}}
\cs{isCentralProje}{Lemma~\ref{l:central}}
\cs{reflectionSummule}{Lemma~\ref{l:refl}}
\cs{existscommonheckerallis}{Lemma~\ref{l:common}}
\cs{hasParahoricTypechosenGLRep}{Lemma~\ref{l:type}}
\cs{ezero}{Definition~\ref{d:good}}
\cs{heckeAlgebra}{Definition~\ref{d:rigid}}
\cs{heckeOperator}{Definition~\ref{d:rigid}}
\cs{eigenPart}{Definition~\ref{d:rigid}}
\cs{occurringLabels}{Definition~\ref{d:rigid}}
\cs{heckeOnGood}{Definition~\ref{d:rigid}}
\cs{traceComparison}{Input 1}
\cs{isCuspKerneltypeKernel}{Input 2}
\cs{finitelevelSpace}{Input 2}
\cs{existsframegradedGram}{Input 3}
\cs{isUnimodulargradedLatticeBasisgram}{Input 3}
\cs{existsorthonormal}{Input 3}
\cs{isCentralProjsumtupleIdem}{Input 4}
\cs{reflectionSummullabelIdem}{Input 4}
\cs{piTensorElemapplyeqzero}{Input 4}
\cs{charpolyeqmuloftracepowadd}{Input 4}
\cs{commuteprep}{Input 5}
\cs{iSupeigenPart}{Input 5}
\cs{finiteeigenPart}{Input 5}
\cs{existsheckeproj}{Input 5}
\cs{occurringLabelssubsingleton}{Input 6}
\cs{auxLift}{Input 7}
\cs{auxLiftinjective}{Input 7}
\cs{liftsauxLift}{Input 7}
\cs{auxLiftprincipalRep}{Input 7}
\cs{iotaAtLevel}{Input 7}
\cs{iotaAtLevelinjective}{Input 7}
\cs{existsisSphericalWithiota}{Input 7}
\cs{algFormRepiota}{Input 7}
\cs{smulchosenGLRepeq}{Input 7}
\cs{hasGapgridParam}{Input 8}
\cs{cardneggridSigns}{Input 8}
\cs{omegazpowneone}{Input 8}
\newcommand{\cV}{{\mathcal V}}
\newcommand{\fs}{{\mathfrak s}}
\newcommand{\sG}{{\mathsf G}}
\newcommand{\cH}{{\mathcal H}}
\newcommand{\Occ}{\mathrm{Occ}}
\newcommand{\Ch}{\mathrm{Ch}}
\newcommand{\GL}{\mathrm{GL}}
\newcommand{\num}{\mathrm{num}}

\begin{document}
\color{ink}

{\raggedright
{\sffamily\bfseries\fontsize{7.2}{9}\selectfont\color{leancol}\addfontfeatures{LetterSpace=8}READBACK NOTE 26\par}
\vspace{5pt}
{\fontsize{21.5}{25.5}\selectfont\bfseries Good forms and rows: equal Hecke traces on unitary groups that differ at the real places\par}
\vspace{5pt}
{\fontsize{12.4}{16}\selectfont\itshape\color{inksoft} One unit of the Lean artifact \ArtifactName, read back into
mathematics\par}}
\vspace{8pt}
{\color{rulegrey}\hrule height 0.6pt}
\vspace{7pt}
\thispagestyle{empty}

\textbf{Summary.} The rows of the proof are cusp forms on an indefinite unitary group, cut out of the forms of
principal level at the chosen primes by one idempotent of the group algebra of a finite group: the \emph{good
forms}. This note defines these objects: the hermitian spaces \m{\spaceY{D}{i}} and \m{\spaceW{D}{i}} attached to the data~$D$, the forms of
principal level on them, the representation of the finite group, the good forms, the rows and their dimension.
Its theorem is a transfer between unitary groups that differ only at the real places: if on two such spaces every
Hecke eigensystem meets at most one primitive label at each chosen prime, then every Hecke operator away from
the chosen primes has the same trace, hence the same determinant, on the good forms of the two spaces. The
hypothesis holds for \m{\spaceY{D}{0}} and for every \m{\spaceW{D}{i}}, so the rows have the dimension of the good
forms on each \m{\spaceW{D}{i}}.
The proof subtracts from the good idempotent its cubic twists, which leaves group elements with an irreducible
characteristic polynomial at a chosen prime; for those a comparison of trace formulas is available.

\statement{leancol}{\fontsize{9.6}{13.2}\selectfont\textbf{What this note is.} The artifact \ArtifactName{} is a
body of Lean code: the compact edition of a formal proof, due to L.~Alpöge with the AI system Claude, of the
Riemann hypothesis and of its generalisation to Dirichlet $L$-functions. That headline is a claim. The authors
report that every file compiles and that Lean's kernel accepted the proofs; an independent replay is recorded in
the handout's record of checks when it reports. No mathematician has read the argument through, and nothing here
says that the Riemann hypothesis is proved. This note treats \UnitInNote{} of the \UnitMembers{} declarations of
the unit \emph{Good idempotents, labels and the transfer of traces}; the other declarations, on labels over
finite fields, are the subject of the note \emph{Primitive labels}. Proved in full here, from the inputs: the
transfer of traces with its three corollaries (Theorems 1 and~2), the reduction of Theorem~3 to its inputs, and
the lemmas of §§\ref{s:spaces}--\ref{s:hecke}. Taken for granted and stated as Inputs~1 to~8: a comparison of
trace formulas, finiteness of spaces of cusp forms and facts on lattices and on finite general linear groups from
the artifact's background library; and results of other units of the argument, among them the uniqueness of
primitive labels (unit on Galois representations) and the injectivity and equivariance of theta lifts. Every
statement was compared with its Lean text, argument by argument. The Lean text read is \PinText; Appendix~A is
the concordance of \ConcDecls{} declarations.}

\sect{The statements}{s:intro}
The objects are defined in §\ref{s:setting}. In brief: $D$ is admissible data (Setting~\ref{s:data});
\m{\Prin{D}{\cV}} is the space of forms of principal level at the chosen primes on a hermitian space~$\cV$,
\m{\Good{D}{\cV}} its subspace of good forms, $\cH(D)$ the algebra of Hecke operators away from the chosen primes,
and \m{\Prin{D}{\cV}[p]} the part of \m{\Prin{D}{\cV}} with Hecke eigensystem~$p$. Rigidity of labels is
Definition~\ref{d:rigid}.

\statement{arith}{\phantomsection\label{t:1}\textbf{Theorem 1} (equal traces).\enspace
\emph{In Setting~\ref{s:data}, let $\fs$, $\fs'$, $\cV$, $\cV'$ be as in Setting~\ref{s:two}. Suppose that
\Rigid{D}{\cV} and that \Rigid{D}{\cV'}. Then for every \m{T\in\cH(D)}, the trace of $T$ on \m{\Good{D}{\cV}}
equals the trace of $T$ on \m{\Good{D}{\cV'}}.}}
\leanline{\lean{traceheckeOnGoodeq}}

\statement{arith}{\phantomsection\label{t:2}\textbf{Theorem 2} (determinants, eigensystems, dimensions).\enspace
\emph{In Setting~\ref{s:data}, let $\fs$, $\fs'$, $\cV$, $\cV'$ be as in Setting~\ref{s:two}, and suppose that
\Rigid{D}{\cV} and that \Rigid{D}{\cV'}.}
\begin{enumerate}[label=\textup{(\alph*)},leftmargin=2.3em,itemsep=2pt]
\item \emph{For every \m{T\in\cH(D)}, the determinant of $T$ on \m{\Good{D}{\cV}} equals the determinant of $T$ on
\m{\Good{D}{\cV'}}.}
\item \emph{Let $p$ be an eigensystem. If \m{\Good{D}{\cV'}\cap\Prin{D}{\cV'}[p]=0}, then
\m{\Good{D}{\cV}\cap\Prin{D}{\cV}[p]=0}.}
\end{enumerate}
\emph{Moreover, in Setting~\ref{s:data}, for every integer $i$ with \m{1\le i\le m(D)}:}
\[\rowDim{D}=\dim_{\mathbb C}\Good{D}{\spaceW{D}{i}} .\]}
\leanline{\lean{detheckeOnGoodeq}\sep \lean{goodFormsinfeigenParteqbot}\sep \lean{finrankrowseqfinrankgoodForms}}

\statement{arith}{\phantomsection\label{t:3}\textbf{Theorem 3} (rigid labels).\enspace
\emph{In Setting~\ref{s:data}: \Rigid{D}{\spaceY{D}{0}}; and for every integer $i$ with \m{1\le i\le m(D)},
\Rigid{D}{\spaceW{D}{i}}.}}
\leanline{\lean{hasRigidLabelsspaceYzero}\sep \lean{hasRigidLabelsspaceW}}

The space \m{\spaceY{D}{0}} carries the rows, \m{\Row{D}=\Good{D}{\spaceY{D}{0}}}; it has $m(D)/2$ negative signs at
each infinite place. The space \m{\spaceW{D}{i}} is the source of the $i$-th single theta lift and has \m{m(D)-i}
negative signs. Both have rank~$r(D)$ and the same finite adelic group, so Theorems 1 and~2 compare them. In
Chapter~VII the equality of determinants turns the volume of the chain of $m(D)$ lifts on the rows into the product
of the volumes of the single lifts, and the equality of dimensions lets one count rows on~\m{\spaceW{D}{i}}.

\sect{The setting}{s:setting}
\begin{conventions}\phantomsection\label{c:conventions}
The natural numbers are \m{0,1,2,\dots}; differences \m{n-n'} of natural numbers that the formal text truncates at
zero are marked as such. The parameters \m{P=(N_0,d)} of the proof are fixed: $N_0$ is a positive multiple
of~$18480$ and $d$ a prime with \m{d\equiv N_0/2-1} modulo~$N_0$. For a number field~$E$ with ring of integers
$O_E$: a \emph{prime} of~$E$ is a non-zero prime ideal, also called a finite place; $\F_v$ is the residue field
of a prime~$v$, $E_v$ the completion and $O_v$ its ring of integers; $\mathbb A_E$ and $\mathbb A_{E,f}$ are the
adeles and the finite adeles, and $\hat O_E$ is the ring of integral finite adeles. For a CM field $E$,
\m{z\mapsto\bar z} is complex conjugation, on $E$, on its adeles and on its primes; a prime $v$ is \emph{split} if
\m{\bar v\ne v}. For a finite set $I$ and a commutative ring, \m{\GL_I} of that ring is the group of the invertible matrices
over it with rows and columns indexed by~$I$, and $g^*$ is the conjugate transpose of a matrix $g$ over a ring
with conjugation. \m{\mathbb C[G]} is the complex group algebra of a group~$G$, and a complex representation
of~$G$ is also a module over it. A \emph{Boolean} is one of the two truth values. Labels, primitive labels and
the idempotents \m{e_{I,c}} and \m{e^{\otimes}_{I,\mathbf c}} of \m{\mathbb C[\GL_I(\F)]} and of the group algebra of
a product, for a finite field~$\F$, are those of the note \emph{Primitive labels}, and its sets of good
tuples are written as there. Letters of the handout's overview that are not displayed with the data~$D$ as an
argument (among them $t$, $\Upsilon$, $R$, $f$, $\beta$, $\nu$ and~$n$) have, in this note, only the meaning given
in its notation list.
\end{conventions}

\begin{setting}[the data]\label{s:data}
$D$ is admissible data for the parameters~$P$, in the sense of Chapter~III of the handout (Setting~1 of its
overview). The following objects attached to $D$ are defined in Chapter~III and in the unit on the seesaw; they
are recalled here and used throughout.
\begin{enumerate}[label=\textup{(\roman*)},leftmargin=2.3em,itemsep=2pt]
\item The CM field \m{E=E(D)}, abelian over~$\mathbb Q$; it contains a primitive root of unity of order~$5040$ and
has an even number of infinite places. The letter $\tau$ denotes an infinite place of~$E(D)$.
\item Natural numbers $m(D)$, a power of~$2$, and $r(D)$, a power of~$3$ greater than~$1$, with
\m{\sqrt{x(D)}\le m(D)\le2\sqrt{x(D)}} and \m{x(D)^5\le r(D)\le3\,x(D)^5}, where the first scale $x(D)$ of the data is
larger than every absolute constant that occurs in this note; the exponent $\mu(D)$ of the data, with
\m{m(D)^{22}=2^{\mu(D)}}; the numbers
\m{k(D)=r(D)+1-m(D)\ge1}, \m{t(D)=5r(D)}, \m{\varrho(D)=m(D)/2-1}, \m{h_2(D)=r(D)^8} and \m{h(D)=2r(D)^8}; the ranks
\m{n_i(D)=k(D)+i\,(t(D)+1)} for \m{i\ge0}; and the integer \m{M_\omega(D)=(N_0d)^{\mu(D)+6}}.
\item For every natural number $i$ and every $\tau$, the \emph{weight} of Chapter~III, the natural number
\[
w_{D,i}(\tau)=N_1(D)\,\bigl\lceil w_D(\tau)/N_1(D)\bigr\rceil+i\,h(D)+h_2(D).
\]
Here \m{w_D(\tau)} is the nominal weight of the overview and \m{N_1(D)} the order of the root of unity that
generates $E(D)$ over \m{\mathbb Q(\sqrt{-d})}; the first term is the least multiple of \m{N_1(D)} that is at least
the nominal weight. By the definition, \m{w_{D,i'}(\tau)-w_{D,i}(\tau)=(i'-i)\,h(D)}; and
\m{w_{D,i}(\tau)\ge h(D)}, which is not read off the formula: it is one of the comparisons of magnitudes of the
data (Chapter~III).
\item The level, a prime number $N(D)$, with an integer $\pi_N(D)$ of~$E(D)$ such that
\m{\pi_N(D)\,\overline{\pi_N(D)}=N(D)}.
\item Two distinct primes $g_0(D)$, $g_1(D)$ and an integer $\pi_g(D)$ of \m{K=\mathbb Q(\sqrt{-d})} with
\m{\pi_g(D)\overline{\pi_g(D)}=g_0(D)g_1(D)}; all three are components of~$D$. The \emph{chosen primes} are the primes of $E(D)$ that divide~$\overline{\pi_g(D)}$; their set is
\m{\Ch(D)}. No chosen prime is the conjugate of a chosen prime; in particular each is split. A chosen prime $v$
divides exactly one of $g_0(D)$, $g_1(D)$, and $\F_v$ has $g_0(D)$ or $g_1(D)$ elements accordingly. The number
\m{\#\F_v-1} is divisible by~$9$ and by \m{M_\omega(D)^2}, and \m{M_\omega(D)+1<\#\F_v}.
\item The \emph{row index set} $I(D)$, a finite set with $r(D)$ elements, with a bijection $\num_D$ from $I(D)$
to \m{\{0,\dots,r(D)-1\}}; and the finite group
\[\sG(D)=\prod_{v\in\Ch(D)}\GL_{I(D)}(\F_v).\]
\item For every prime $v$ of $E(D)$, a multiplicative character $\omega_{D,v}$ of~$\F_v$: of order
\m{M_\omega(D)} if $v$ is chosen, trivial otherwise.
\item The \emph{row parameter}: for an infinite place $\tau$ and a natural number~$n$, the integer
\begin{align*}
\varkappa_D(\tau,n)&=\lfloor k(D)/2\rfloor-1-n&&\text{if }n+1<k(D),\\
\varkappa_D(\tau,n)&=h_2(D)-w_{D,\,n+2-k(D)}(\tau)&&\text{otherwise.}
\end{align*}
\item \emph{Tame characters.} A tame character $\vartheta$ of~$D$ is a continuous character of the norm-one ideles
of~$E(D)$, trivial on the norm-one elements of~$E(D)$, together with a finite set $Q(\vartheta)$ of primes
of~$E(D)$ and a family \m{(\psi_{\vartheta,v})_v} of multiplicative characters of the residue fields, extended
by~$0$, its \emph{tame part}. The character and the family are tied as in the note on the auxiliary pair:
$\vartheta$ has weight~$0$ and tame part \m{(\psi_{\vartheta,v})_v}, that is,
\m{\vartheta(z)=\prod_v\psi_{\vartheta,v}(z\bmod v)} for every norm-one idele~$z$ whose finite part is integral; in
particular $\vartheta$ is trivial at the infinite places. It is required that \m{\psi_{\vartheta,v}=1} for
\m{v\notin Q(\vartheta)}, that no conjugate of a member of $Q(\vartheta)$ lies in $Q(\vartheta)$, and that the
product of the norms of the members of $Q(\vartheta)$ is prime to \m{g_0(D)g_1(D)N(D)}. The trivial character,
with \m{Q=\emptyset}, is one.
\end{enumerate}
\end{setting}

\subhead{Hermitian spaces with a point, and forms.} These notions belong to the background library; they are
recalled as the library defines them. Let $E$ be a CM field, $I$ a finite set, and \m{\fs\colon I\to\{1,-1\}} a
family of signs. Put
\begin{align*}
G_\infty(\fs)&=\prod_\tau\mathrm U(\fs),\\
G_f(I)&=\{g\in\GL_I(\mathbb A_{E,f}):g^*g=1\},\qquad K_0(I)=G_f(I)\cap\GL_I(\hat O_E),
\end{align*}
the product over the infinite places of~$E$, where \m{\mathrm U(\fs)} is the group of the complex matrices $g$
with \m{g^*\,\mathrm{diag}(\fs)\,g=\mathrm{diag}(\fs)}. For a prime $v$ and \m{g\in G_f(I)}, \m{g_v\in\GL_I(E_v)} is
the component at~$v$; an element of $K_0(I)$ has a reduction in \m{\GL_I(\F_v)}. For a split prime $v$ and
\m{k\in\GL_I(E_v)}, the element \m{k^{(v)}\in G_f(I)} has component $k$ at~$v$ and $1$ at the primes other than
$v$ and~$\bar v$.

A \emph{space over $E$ with index set $I$ and signs~$\fs$} is a pair \m{\cV=(J,\mathbf x)} of a matrix $J$
over~$E$ and an element \m{\mathbf x\in\GL_I(\mathbb A_E)}, both indexed by~$I$, such that
\m{\mathbf x^*J\mathbf x} is \m{\mathrm{diag}(\fs)} at every infinite place and~$1$ at the finite places. So $J$
is a hermitian form whose signature at each infinite place is that of~$\fs$, $\mathbf x$ identifies its unitary
group over the adeles with \m{G_\infty(\fs)\times G_f(I)}, and the rational points become the discrete subgroup
\m{\Gamma_\cV=\mathbf x^{-1}\mathrm U(J)(E)\,\mathbf x}. The space is \emph{standard} if all signs are~$1$ and
\m{\mathbf x=1}. The \emph{diagonal space} of the signs $\fs$ has \m{J=\mathrm{diag}(\fs)} and a diagonal
point, with entry~$1$ where the sign is~$1$ and, where it is~$-1$, a fixed idele that is $1$ at the infinite
places and whose product with its conjugate is~$-1$ at the finite places.

Let \m{\epsilon\colon G_\infty(\fs)\to\mathbb C} be a function (a \emph{type}). A \emph{form of type $\epsilon$
on~$\cV$} is a bounded continuous function \m{F\colon G_\infty(\fs)\times G_f(I)\to\mathbb C} that is invariant
under left translation by~$\Gamma_\cV$ and under right translation by an open subgroup of~$G_f(I)$, and that
satisfies \m{\int_{G_\infty(\fs)}\epsilon(u)\,F(g\cdot(u,1))\,du=F(g)} for all~$g$ if $\epsilon$ is integrable. The
group $G_f(I)$ acts on these forms by right translation, \m{(R(g)F)(g')=F(g'g)}, and they carry the Petersson
inner product \m{\langle F,F'\rangle=\int_{\Gamma_\cV\backslash G}\overline{F}F'}, for which $R$ is unitary. Note
that $G_f(I)$ does not depend on the space: it acts on the forms of every space with index set~$I$.

For an integer vector \m{l=(l_\iota)_{\iota\in I}}, the library defines a function $\epsilon_l$ on
\m{\mathrm U(\fs)}: the formal degree times the complex conjugate of the canonical matrix coefficient of the
discrete series with Harish-Chandra parameter~$l/2$. The vector $l$ \emph{has gap $H$ for~$\fs$}, for a real
number~$H$, if its entries are distinct, every \m{l_\iota+\#I-1} is even, and \m{|l_\iota-l_{\iota'}|\ge2H}
whenever \m{\fs_\iota\ne\fs_{\iota'}}. A function $\epsilon$ on \m{G_\infty(\fs)} is a \emph{cusp kernel of
order~$B$} if \m{\|u\|^{n}\epsilon(u)} is integrable for \m{n\le B+1}, \m{\|u\|^{B}\epsilon(u)} is bounded, and
\m{\alpha*\epsilon=\epsilon} for some continuous function $\alpha$ of compact support; here \m{\|u\|} is the
product of the operator norms of the components of~$u$.

In Setting~\ref{s:data}, for signs $\fs$ on $I(D)$, the \emph{type of the data} is the function on
\m{G_\infty(\fs)}
\[
\epsilon_{D,\fs}(u)=\prod_\tau\epsilon_{l_D(\tau)}(u_\tau),\qquad l_D(\tau)=\bigl(2\,\varkappa_D(\tau,\num_D(\iota))\bigr)_{\iota\in I(D)} .
\]
The \emph{principal congruence subgroup at the chosen primes}, \m{K_g(D)\subset K_0(I(D))}, consists of the
elements whose reduction at every chosen prime is~$1$. For \m{\gamma\in\sG(D)}, \m{\tilde\gamma\in K_0(I(D))} is
a lift that the library chooses: its reductions at the chosen primes are the components of~$\gamma$, and its
components at the primes that are neither chosen nor conjugate to a chosen prime are~$1$.

\begin{setting}[two spaces]\label{s:two}
In Setting~\ref{s:data}: $\fs$ and $\fs'$ are families of signs on~$I(D)$, $\cV$ is a space over $E(D)$ with
index set $I(D)$ and signs~$\fs$, and $\cV'$ one with signs~$\fs'$; and for every infinite place~$\tau$ the
vector $l_D(\tau)$ has gap $h(D)$ for~$\fs$ and has gap $h(D)$ for~$\fs'$.
\end{setting}

\subhead{The spaces.}
\begin{definition}[index sets and their bijections]\label{d:index}
\textup{(a)} A \emph{shape} is a triple \m{\Sigma=(\varrho,\mathcal X,o)} of a natural number, a finite set and
an element of that set; its index set is the disjoint union
\m{\{\mathsf b_0,\dots,\mathsf b_{\varrho-1}\}\sqcup\mathcal X\sqcup\{z\}\sqcup\{\mathsf b^*_0,\dots,\mathsf b^*_{\varrho-1}\}}. The bijection
\m{\besselIx{\Sigma}} from \m{\mathcal X\sqcup\{0,\dots,2\varrho\}} to this index set is the identity on
$\mathcal X$ and sends $n$ to $\mathsf b_n$ for \m{n<\varrho}, $\varrho$ to~$z$, and $n$ to \m{\mathsf b^*_{2\varrho-n}} for
\m{n>\varrho}.

\textup{(b)} In Setting~\ref{s:data}, let $i$ be a natural number. $\mathcal X_i(D)$ is a set with $n_i(D)$
elements, one of them marked as the last, and \m{I_i(D)=I(D)\sqcup(\{1,\dots,i\}\times\{0,\dots,t(D)-1\})}; the row
index set $I(D)=I_0(D)$ is the index set of the shape \m{\Sigma(D)=(\varrho(D),\mathcal X_0(D),\text{last})}.
Suppose \m{i<m(D)}. Then \m{\indexY{D}{i}{i<m(D)}} is a bijection from
\m{\mathcal X_i(D)\sqcup\{0,\dots,2\varrho(D)-i\}} to~$I_i(D)$: for \m{i=0} it is \m{\besselIx{\Sigma(D)}}, and for
\m{i\ge1} it is a bijection that the formal text chooses, the two sets having the same number of elements.
\end{definition}
\leanline{\lean{besselIndexEquiv}\sep \lean{indexEquivY}}

\begin{lemma}\label{l:m}
In Setting~\ref{s:data}: \m{m(D)=2\varrho(D)+2}; and
\[
\exp(-4)\,\sqrt{x(D)}\le\varrho(D)\le\exp(2)\,\sqrt{x(D)}.
\]
\end{lemma}
\begin{proof}
The number $m(D)$ is even and at least~$2$, and \m{\varrho(D)=m(D)/2-1}. Now $x(D)$ is large and
\[
\sqrt{x(D)}\le m(D)\le2\sqrt{x(D)},
\]
so \m{m(D)/4\le\varrho(D)\le m(D)}, which gives the bounds.
\end{proof}
\leanline{\lean{meqtwomulvarrhoaddtwo}\sep the bounds are an instance without a name, in the form in which
the formal text compares magnitudes of the data}

\begin{definition}[ideles]\label{d:ideles}
In Setting~\ref{s:data}:
\begin{enumerate}[label=\textup{(\alph*)},leftmargin=2.3em,itemsep=2pt]
\item \m{\chosenIdele{D}} is a finite idele of $E(D)$, chosen by the library, that is integral, has component~$1$
at every prime not dividing~$\overline{\pi_g(D)}$, and generates the ideal \m{\overline{\pi_g(D)}\hat O_E}.
\item \m{\ideleA{D}=\alpha(D)\big/\overline{\chosenIdele{D}}}, where $\alpha(D)$ is the finite idele of the Weil
data of $E(D)$ (Chapter~III).
\item \m{\piIdele{D}} is the principal finite idele of the element $\pi_N(D)$ of~$E(D)$.
\item For a natural number~$i$, \m{\latticeScale{D}{i}=\bigl(\ideleA{D}\big/\overline{\ideleA{D}}\bigr)^{\lfloor(m(D)-i+1)/2\rfloor}},
with \m{m(D)-i} truncated at zero: a finite idele of norm one, equal to~$1$ for \m{i\ge m(D)}.
\end{enumerate}
\end{definition}
\leanline{\lean{chosenIdele}\sep \lean{ideleA}\sep \lean{piIdele}\sep \lean{latticeScale}}

\begin{lemma}\label{l:pi}
In Setting~\ref{s:data}: \m{N(D)=\piIdele{D}\cdot\overline{\piIdele{D}}} as finite ideles; for every prime $v$
of~$E(D)$ the component of \m{\piIdele{D}} at $v$ is a uniformizer if $v$ divides $\pi_N(D)$ and a unit otherwise;
and if $v$ divides $\pi_N(D)$, then $\bar v$ does not.
\end{lemma}
\begin{proof}
The first assertion is \m{\pi_N(D)\overline{\pi_N(D)}=N(D)}. The component of \m{\piIdele{D}} at $v$ is the
image of~$\pi_N(D)$, so the second and the third assertion are Input~8(e). (They hold because the prime $N(D)$ is
unramified in~$E(D)$: if $v^2$ or $v\bar v$ divided $\pi_N(D)$, then $v^2$ would divide~$N(D)$.)
\end{proof}
\leanline{\lean{piIdelespec}}

Two matrices of the background are needed to state the next lemma. Let $\beta$ be a bijection from
\m{\mathcal X\sqcup\{0,\dots,s-1\}} to a finite set~$I'$, where $\mathcal X$ is a finite set with a marked last
element and $s$ a natural number. The \emph{graded Gram matrix}
\m{\mathrm{Gr}(\beta,N)}, for $N$ in a commutative ring, is the matrix indexed by~$I'$ that is the identity on
\m{\beta(\mathcal X)}, pairs \m{\beta(n)} with \m{\beta(n')} to~$1$ if \m{n+n'=s-1} and \m{n\ne n'}, to $-N$ if
\m{n=n'} and \m{2n=s-1}, and has all other entries~$0$. The \emph{graded lattice basis}
\m{\mathrm{Lb}(\beta,a,\pi,u)}, for units $a$, $\pi$ and an element $u$ of norm one of a ring with conjugation, is
the matrix indexed by~$I'$ that for even $s$ is diagonal, with entry $1$ on \m{\beta(\mathcal X)}, \m{u\,a^{-1}} at
\m{\beta(n)} for \m{2n<s} and \m{u\,\bar a} for \m{2n\ge s}; for odd $s$ it is diagonal with entry $1$ on
\m{\beta(\mathcal X)} except \m{\bar\pi^{-1}} at the marked element, and $u$ on the \m{\beta(n)}, plus the single
entry \m{u\,(\pi\bar\pi)^{-1}} in the row of the middle index \m{\beta((s-1)/2)} and the column of the marked
element. The signs \m{\fs^{Y}_{D,i}} on $I_i(D)$ are $-1$ at the \m{\iota\in I(D)} for which
\m{\num_D(\iota)-k(D)-i} is non-negative and even, and~$1$ elsewhere; the signs \m{\fs^{W}_{D,i}} on $I(D)$ are
$-1$ at the $\iota$ with \m{\num_D(\iota)\ge k(D)+i-1} (truncated), and $1$ elsewhere.

\begin{lemma}[existence of the spaces of the chain]\label{l:Y}
In Setting~\ref{s:data}, let $i$ be a natural number. There is a space $\cV$ over $E(D)$ with index set
$I_i(D)$ and signs \m{\fs^{Y}_{D,i}} such that:
\begin{enumerate}[label=\textup{(\alph*)},leftmargin=2.3em,itemsep=2pt]
\item if \m{m(D)\le i}, then $\cV$ is standard;
\item if \m{i<m(D)}, then the matrix of $\cV$ is \m{\mathrm{Gr}(\indexY{D}{i}{i<m(D)},N(D))}, and there is
\m{k\in\GL_{I_i(D)}(\hat O_E)} whose reduction at every chosen prime is~$1$ such that the finite part of the point
of $\cV$ is \m{\mathrm{Lb}\bigl(\indexY{D}{i}{i<m(D)},\ideleA{D},\piIdele{D},\latticeScale{D}{i}\bigr)\cdot k}.
\end{enumerate}
\end{lemma}
\leanline{\lean{existsspaceY}}

The proof is given in §\ref{s:spaces}.

\begin{definition}[the two families of spaces]\label{d:YW}
In Setting~\ref{s:data}, let $i$ be a natural number. \m{\spaceY{D}{i}} is a space as in Lemma~\ref{l:Y}, chosen
once and for all. \m{\spaceW{D}{i}} is the diagonal space over $E(D)$ of the signs \m{\fs^{W}_{D,i}} on~$I(D)$.
\end{definition}
\leanline{\lean{spaceY}\sep \lean{spaceW}}

So \m{\spaceY{D}{i}} has rank \m{r(D)+i\,t(D)} and \m{\lfloor(m(D)-i)/2\rfloor} negative signs (truncated), and
\m{\spaceW{D}{i}} has rank $r(D)$ and \m{m(D)-i} negative signs (truncated), at every infinite place.

\begin{lemma}\label{l:std}
In Setting~\ref{s:data}, the spaces \m{\spaceY{D}{m(D)}} and \m{\spaceW{D}{m(D)}} are standard.
\end{lemma}
\begin{proof}
The first is Lemma~\ref{l:Y}(a). For the second, \m{k(D)+m(D)-1=r(D)}, so no \m{\iota\in I(D)} has
\m{\num_D(\iota)\ge k(D)+m(D)-1}: all signs are~$1$, and the diagonal space of these signs is standard.
\end{proof}
\leanline{\lean{isStdspaceY}\sep for \m{\spaceW{D}{m(D)}}, an instance without a name}

\subhead{Forms of principal level, good forms, rows.}
\begin{definition}[principal forms and the action of $\sG(D)$]\label{d:prin}
In Setting~\ref{s:data}, let $\fs$ be a family of signs on $I(D)$ and $\cV$ a space over $E(D)$ with index set
$I(D)$ and signs~$\fs$.
\begin{enumerate}[label=\textup{(\alph*)},leftmargin=2.3em,itemsep=2pt]
\item \m{\Prin{D}{\cV}} is the space of the forms of type \m{\epsilon_{D,\fs}} on~$\cV$ that are fixed by
\m{R(k)} for all \m{k\in K_g(D)}.
\item \m{\prep{D}{\cV}} is the representation of $\sG(D)$ on \m{\Prin{D}{\cV}} given by
\m{\prep{D}{\cV}(\gamma)=R(\tilde\gamma)}.
\end{enumerate}
\end{definition}
\leanline{\lean{principalForms}\sep \lean{principalRep}}

\begin{definition}[the good idempotent, good forms, rows]\label{d:good}
In Setting~\ref{s:data}, the \emph{good idempotent} is the element of \m{\mathbb C[\sG(D)]}
\[
e_0(D)=\sum_{\mathbf c}e^{\otimes}_{I(D),\mathbf c},
\]
the sum over the families \m{\mathbf c=(c_v)_{v\in\Ch(D)}} of monic polynomials \m{c_v\in\F_v[X]} of degree~$r(D)$ such
that, for each of the two primes $g_0(D)$ and~$g_1(D)$, the family of the $c_v$ with $v$ above that prime belongs to
\m{\mathrm{Good}_{M_\omega(D),r(D)}}. (This definition belongs to the unit on the seesaw.) Let $\fs$ be a family
of signs on $I(D)$ and $\cV$ a space over $E(D)$ with index set $I(D)$ and signs~$\fs$. Then
\[
\Good{D}{\cV}=\prep{D}{\cV}\bigl(e_0(D)\bigr)\,\Prin{D}{\cV},
\]
the space of \emph{good forms} on~$\cV$. The space of \emph{rows} is \m{\Row{D}=\Good{D}{\spaceY{D}{0}}}, and
\m{\rowDim{D}} is its complex dimension.
\end{definition}
\leanline{\lean{goodForms}\sep \lean{rows}\sep \lean{rowDim}\sep the good idempotent: \lean{ezero}}

So a good family has, above each of $g_0(D)$ and~$g_1(D)$, a label without a root of order dividing $M_\omega(D)$ at
every chosen prime and a primitive label at one chosen prime at least. The rows of the overview of this handout
are \m{\Row{D}}.

\subhead{Hecke operators and labels (definitions of the unit on fit twists).}
In Setting~\ref{s:data}, $\cH(D)$ is the algebra of the families \m{T=(T_\cV)}, indexed by all spaces $\cV$ over
$E(D)$ with index set $I(D)$ and any signs, of endomorphisms $T_\cV$ of \m{\Prin{D}{\cV}} for which there is one
element \m{\sum_gc_g\,g} of the group algebra of the subgroup of the \m{g\in G_f(I(D))} with \m{g_v=1} at every
chosen prime~$v$, such that \m{T_\cV=\sum_gc_g\,R(g)} on \m{\Prin{D}{\cV}} for every~$\cV$. Every $T_\cV$
commutes with \m{\prep{D}{\cV}}, so it preserves \m{\Good{D}{\cV}}; the trace and the determinant of $T$ on
\m{\Good{D}{\cV}} are those of this restriction. For a prime $v$ of~$E(D)$ that is split, with neither $v$ nor
$\bar v$ chosen, and a natural number~$n$, the Hecke operator \m{\mathbb T^{v,n}\in\cH(D)} is the operator of the double
coset of \m{\mathrm{diag}(\varpi,\dots,\varpi,1,\dots,1)} with $n$ entries~$\varpi$, a uniformizer, in
\m{\GL_{I(D)}(E_v)} acting through \m{k\mapsto k^{(v)}} (it is zero if $n$ exceeds the rank); for the other pairs
\m{(v,n)} put \m{\mathbb T^{v,n}=0}. An
\emph{eigensystem} is a complex function $p$ on the pairs \m{(v,n)} of a prime of $E(D)$ and a natural number,
and
\[
\Prin{D}{\cV}[p]=\{y\in\Prin{D}{\cV}:\ \mathbb T^{v,n}_\cV\,y=p(v,n)\,y\ \text{for all }(v,n)\}.
\]
For an eigensystem $p$ and a chosen prime~$v$, \m{\Occ(D,\cV,p,v)} is the set of the polynomials \m{c\in\F_v[X]}
that are primitive labels of degree~$r(D)$ and for which \m{e_{I(D),c}}, acting through \m{\prep{D}{\cV}}
restricted to the factor \m{\GL_{I(D)}(\F_v)} of~$\sG(D)$, does not kill \m{\Prin{D}{\cV}[p]}.

\begin{definition}[rigid labels]\label{d:rigid}
In Setting~\ref{s:data}, let $\fs$ be a family of signs on $I(D)$ and $\cV$ a space over $E(D)$ with index set
$I(D)$ and signs~$\fs$. Then \Rigid{D}{\cV} if for every eigensystem $p$ and every chosen prime~$v$ the set
\m{\Occ(D,\cV,p,v)} has at most one element.
\end{definition}
\leanline{\lean{HasRigidLabels}\sep the notions it rests on: \lean{heckeAlgebra}\sep \lean{heckeOperator}\sep
\lean{eigenPart}\sep \lean{occurringLabels}\sep \lean{heckeOnGood}}

\sect{Inputs}{s:inputs}
\statement{inksoft}{\textbf{Inputs from the background library.} In Inputs 1 to~3, $E$ is a CM field that
contains a primitive root of unity of order~$5040$, $I$ is a non-empty finite set, and spaces, forms and types are
those recalled in §\ref{s:setting}.

\vspace{4pt}\phantomsection\label{i:1}\textbf{Input 1} (comparison of traces on unitary groups that differ at the
real places).\enspace \emph{Suppose that $E$ has an even number of infinite places. Let $\mathbf K$ be an open subgroup of
$K_0(I)$ and \m{g\in G_f(I)} an element such that every element of the coset $g\mathbf K$ has, at some prime $v$ of~$E$, a
component in \m{\GL_I(E_v)} whose characteristic polynomial is irreducible. Let $B$ be a natural number with
\m{B\ge16\,[E:\mathbb Q]\,(\#I+1)^3}, let \m{l=(l(\tau))_\tau} be a family of integer vectors indexed by~$I$, one for
each infinite place, and let $H$ be a real number with \m{H\ge500\,(\#I)^3+10\,(\#I)^5}. Then there is a complex
number $t$ with the following property. For every family of signs $\fs$ on~$I$, every space $\cV$ over $E$ with
index set $I$ and signs~$\fs$, and every finite orthonormal basis \m{(b_n)_n} of the space of the forms of type
\m{\epsilon=\prod_\tau\epsilon_{l(\tau)}} on $\cV$ that are fixed by~$\mathbf K$: if $\epsilon$ is a cusp kernel of
order~$B$ and every $l(\tau)$ has gap $H$ for~$\fs$, then \m{\sum_n\langle b_n,R(g)\,b_n\rangle=t}.}

\vspace{3pt}\phantomsection\label{i:2}\textbf{Input 2} (cusp kernels; finiteness).\enspace
\emph{\textup{(a)} Let $\fs$ be signs on~$I$, $B$ a natural number and \m{(l(\tau))_\tau} integer vectors such that
every $l(\tau)$ has gap \m{500\,(\#I)^3+B} for~$\fs$. Then \m{\prod_\tau\epsilon_{l(\tau)}} is a cusp kernel of
order~$B$. \textup{(b)} If $\epsilon$ is a cusp kernel of some order on \m{G_\infty(\fs)}, $\cV$ a space with signs
$\fs$ and $\mathbf K$ an open subgroup of~$G_f(I)$, then the space of the forms of type $\epsilon$ on $\cV$ fixed
by $\mathbf K$ has finite dimension.}

}

\statement{inksoft}{\textbf{Inputs from the background library, continued.}

\vspace{4pt}\phantomsection\label{i:3}\textbf{Input 3} (frames and lattices).\enspace \emph{Let $\beta$ be a
bijection from \m{\mathcal X\sqcup\{0,\dots,s-1\}} to a finite set~$I'$, as before Lemma~\ref{l:Y}.
\textup{(a)} Let \m{N>0} be an integer and $\fs$ a family of signs on $I'$ with exactly \m{\lfloor(s+1)/2\rfloor}
signs~$-1$. Then there is a complex matrix $y$ with \m{y^*\,\mathrm{Gr}(\beta,N)\,y=\mathrm{diag}(\fs)}.
\textup{(b)} Let $a$, $\pi$ be finite ideles of $E$, $u$ a finite idele of norm one, and $N$ an integer with
\m{N=\pi\bar\pi}. Then for \m{B=\mathrm{Lb}(\beta,a,\pi,u)} the matrix \m{B^*\,\mathrm{Gr}(\beta,N)\,B} has entries
in $\hat O_E$ and an inverse with entries in~$\hat O_E$. \textup{(c)} Every hermitian matrix $b$ over $\hat O_E$
with an inverse over $\hat O_E$ has the form \m{k^*bk=1} for some \m{k\in\GL_{I'}(\hat O_E)}.}

\vspace{3pt}\phantomsection\label{i:4}\textbf{Input 4} (finite general linear groups).\enspace \emph{Let
\m{(\F_v)_{v}} be a finite family of finite fields and $I$ a finite set. \textup{(a)} For every finite set
$\mathcal C$ of families \m{\mathbf c=(c_v)_{v}} of polynomials, \m{\sum_{\mathbf c\in\mathcal C}e^{\otimes}_{I,\mathbf c}}
is an idempotent of the group algebra of \m{\prod_v\GL_I(\F_v)} whose coefficients are constant on conjugacy
classes. \textup{(b)} Let $\chi$ be a multiplicative character of a finite field~$\F$, \m{x_\chi\in\F^\times} the
unit that names it, \m{c\in\F[X]} a polynomial of which $x_\chi$ is not a root, and \m{l\in\F^I} a non-zero vector.
Then \m{\bigl(\sum_h\chi(\det h)^{-1}h\bigr)\cdot e_{I,c}=0}, the sum over the \m{h\in\GL_I(\F)} of the form
\m{h=1+u\,l^{\mathsf T}} with \m{u\in\F^I}. \textup{(c)} If \m{\rho} is a representation of \m{\prod_v\GL_I(\F_v)},
\m{y_v\in\mathbb C[\GL_I(\F_v)]} for all~$v$, and for one index $v_0$ the element $y_{v_0}$, acting through the
factor $v_0$, kills a vector, then \m{\bigotimes_vy_v} kills that vector.
\textup{(d)} Over an algebraically closed field of characteristic~$0$, two square matrices $S$, $S'$, possibly of
different sizes, with \m{\mathrm{tr}(S^n)=\mathrm{tr}(S'^n)} for all \m{n\ge0} have the same characteristic
polynomial.}}
\leanline{\lean{traceComparison}\sep \lean{isCuspKerneltypeKernel}\sep \lean{finitelevelSpace}\sep
\lean{existsframegradedGram}\sep \lean{isUnimodulargradedLatticeBasisgram}\sep \lean{existsorthonormal}\sep
\lean{isCentralProjsumtupleIdem}\sep \lean{reflectionSummullabelIdem}\sep \lean{piTensorElemapplyeqzero}\sep
\lean{charpolyeqmuloftracepowadd}}

\subhead{On Input 1.} It is the one deep input of this note, and its relation to the literature needs care. The
library's file names as sources J.-P.~Labesse, \emph{Changement de base CM et séries discrètes} (2011),
Théorème~5.4 and its proof, and J.~Arthur and L.~Clozel, \emph{Simple algebras, base change, and the advanced theory
of the trace formula} (1989), Chapter~1, Lemmas 2.4 to~2.6; the catalogue of this handout first classed Input~1 as
a special case of these. The second reading of this write-up compared the statements and left the identification
open: Labesse's Théorème~5.1 is an identity of stabilized trace formulas and his Théorème~5.4 gives weak cuspidal
base change and multiplicity one, with pseudo-coefficients as test functions; the lemmas of Arthur and Clozel are
simple trace formulas for general linear groups. None of them states the equality of this trace for a fixed coset
$g\mathbf K$, for integrable kernels with an explicit gap, and for all spaces at once. According to the headers of the library's two files on this comparison (I read the statement of Input~1 and
not its proof), the statement is proved there by a comparison of simple trace formulas in which the test function
is elliptic at a finite place, because of the irreducible characteristic polynomial, and is at the real places
the integrable idempotent kernel of a very regular discrete series in place of a pseudo-coefficient; the rational
elliptic classes of the two groups, which differ at the real places and are therefore not the same sets, are
matched by changing their real signs with a global element that is locally a norm; and convergence is obtained by
packing, without truncation and without endoscopy. So Input~1 is a statement of the formal library
in the direction of known results, not yet a quotation of a published theorem. Input~2(b) is a special case of the
finiteness of spaces of automorphic forms: the second reading places it at Godement's finite-dimensionality
lemma in A.~Borel, \emph{Automorphic forms on reductive groups} (the Park City lectures of the catalogue), §8.3
with §§7.6--7.9, which applies once the library's uniform bound for the height and its compactness criterion are
granted; Harish-Chandra's theorem for a fixed type in §8.4 has different hypotheses. Input~3(c) is a special case
of the classification of unimodular hermitian lattices: C.~T.~C.~Wall, \emph{On the classification of hermitian
forms~I}, Compositio Math.~22 (1970), pp.~432--433, with Theorem~2 for the lift from the residue field, and
R.~Jacobowitz, \emph{Hermitian forms over local fields}, Amer. J. Math. 84 (1962), for the general
classification; Input~4(a) is a weaker consequence of Green's theorems, as in the note
\emph{Primitive labels}; Inputs 2(a), 3(a), 3(b), 4(b), 4(c), 4(d) are proved in the library without a named
source.

\statement{inksoft}{\textbf{Inputs from other units of the argument.} In Setting~\ref{s:data} throughout.

\vspace{4pt}\phantomsection\label{i:5}\textbf{Input 5} (the Hecke algebra; unit on fit twists).\enspace \emph{Let
$\cV$ be a space over $E(D)$ with index set $I(D)$. \textup{(a)} Every $T_\cV$, \m{T\in\cH(D)}, commutes with
\m{\prep{D}{\cV}(y)} for all \m{y\in\mathbb C[\sG(D)]}; so \m{\prep{D}{\cV}} preserves every \m{\Prin{D}{\cV}[p]}.
\textup{(b)} If \m{\Prin{D}{\cV}} has finite dimension, it is the sum of the subspaces \m{\Prin{D}{\cV}[p]}, and
only finitely many of them are non-zero. \textup{(c)} For every eigensystem $p$ and every finite set
$\mathcal F$ of eigensystems there is \m{\Pi\in\cH(D)} such that, for every space $\cV$ and every
\m{p'\in\mathcal F\cup\{p\}}, $\Pi_\cV$ is the identity on \m{\Prin{D}{\cV}[p']} if \m{p'=p} and zero on it if
\m{p'\ne p}.}

\vspace{3pt}\phantomsection\label{i:6}\textbf{Input 6} (at most one primitive label; unit on Galois
representations).\enspace \emph{Let $\fs$ be signs on $I(D)$ and $\cV$ a space over $E(D)$ with index set $I(D)$
and signs~$\fs$. Let $I'$ be a non-empty finite set, \m{\nu\colon I(D)\to I'} an injection, and $L$ a linear map
from \m{\Prin{D}{\cV}} to a space of algebraic modular forms, of some weight, level and level character, on the
standard space with index set~$I'$. Let $p$ be an eigensystem, \m{(\Upsilon_w)_w} a family of functions
\m{\mathbb C\to\mathbb C} indexed by the primes of~$E(D)$, $v$ a chosen prime, $n_1$ a natural number,
\m{\beta\colon I'\to\{0,\dots,n_1\}} a map, \m{\chi_1,\dots,\chi_{n_1}} and $\psi$ multiplicative characters
of~$\F_v$. Suppose:}
\begin{enumerate}[label=\textup{(\roman*)},leftmargin=2.3em,itemsep=1pt]
\item \emph{$L$ is injective;}
\item \emph{for every \m{y\in\Prin{D}{\cV}[p]} and every split prime $w$ with neither $w$ nor $\bar w$ chosen,
$L(y)$ is fixed by the elements of $K_0(I')$ with component~$1$ off $w$ and~$\bar w$, and for every
\m{X\in\mathbb C} the operator \m{\sum_{n}(-X)^n(\#\F_w)^{n(n-1)/2}\,\mathbb T^{w,n}} sends $L(y)$ to
\m{\Upsilon_w(X)\,L(y)}; here \m{\mathbb T^{w,n}} is the Hecke operator of the double coset of
\m{\mathrm{diag}(\varpi,\dots,\varpi,1,\dots,1)}, with $n$ entries~$\varpi$, in \m{\GL_{I'}(E_w)}, acting by right
translation on the functions on \m{G_f(I')} (note \emph{At most one primitive label for a Hecke eigensystem},
§2(a));}
\item \emph{\m{\beta^{-1}(0)=\nu(I(D))}, and for every \m{y\in\Prin{D}{\cV}}, every \m{x\in\GL_{I(D)}(\F_v)} and
every \m{k\in\GL_{I'}(O_v)} whose reduction $\bar k$ is block upper triangular for $\beta$ with block~$0$ equal
to~$x$:
\m{R(k^{(v)})\,L(y)=\prod_{n=1}^{n_1}\chi_n(\det\bar k_{nn})\cdot\psi(\det x)\cdot L\bigl(\prep{D}{\cV}(x_{(v)})\,y\bigr)},
where $\bar k_{nn}$ is the $n$-th diagonal block of~$\bar k$ and \m{x_{(v)}\in\sG(D)} has component $x$ at $v$
and~$1$ elsewhere;}
\item \emph{the characters \m{\chi_1,\dots,\chi_{n_1}} are distinct, and \m{\psi^{M_\omega(D)}=1} and
\m{\chi_n^{M_\omega(D)}=1} for all~$n$.}
\end{enumerate}
\emph{Then \m{\Occ(D,\cV,p,v)} has at most one element.}

}

\statement{inksoft}{\textbf{Inputs from other units of the argument, continued.} In Setting~\ref{s:data}
throughout.

\vspace{4pt}\phantomsection\label{i:7}\textbf{Input 7} (theta lifts; units on the seesaw, on the exchange and on
the Siegel--Weil formula).\enspace \emph{Write $1$ for the trivial tame character.
\textup{(a)} For \m{1\le i\le m(D)} there is a linear map $\Lambda_i$, a theta lift, from
\m{\Prin{D}{\spaceW{D}{i}}} to the algebraic forms on the standard space with index set
\m{I(D)\sqcup\{0,\dots,t(D)-1\}} that have the weight, level and level character of the pair \m{\auxPair{D}{1}{i}}.
It is injective; it satisfies Input~6\textup{(ii)} for every eigensystem~$p$, with functions $\Upsilon_w$ that depend
on $p$ only; and \m{\Lambda_i(\prep{D}{\spaceW{D}{i}}(\gamma)\,y)=\rho_i(\gamma)\,\Lambda_i(y)} for
\m{\gamma\in\sG(D)}, where $\rho_i$ is the action of $\sG(D)$ on vectors of that level described in
Lemma~\ref{l:type}. \textup{(b)} There is a linear map $\Lambda_{\mathrm{ch}}$, the composite of the $m(D)$ theta lifts for the
trivial twist, from \m{\Prin{D}{\spaceY{D}{0}}} to algebraic forms on the standard space \m{\spaceY{D}{m(D)}}. It is
injective; it satisfies Input~6\textup{(ii)} for every~$p$, with some functions~$\Upsilon_w$; and for every
\m{k\in K_0(I_{m(D)}(D))} whose reduction $\bar k_v$ at every chosen prime $v$ is block upper triangular for the blocks
\m{I(D)}, \m{\{1\}\times\{0,\dots,t(D)-1\}}, \dots, \m{\{m(D)\}\times\{0,\dots,t(D)-1\}} of $I_{m(D)}(D)$, with block~$0$ equal
to~$\gamma_v$,
\m{R(k)\,\Lambda_{\mathrm{ch}}(y)=\prod_{v\in\Ch(D)}\prod_{n=1}^{m(D)}\omega_{D,v}^{\,n}(\det\bar k_{v,nn})\cdot\Lambda_{\mathrm{ch}}\bigl(\prep{D}{\spaceY{D}{0}}(\gamma)\,y\bigr)}.
\textup{(c)} Let $O$ be an auxiliary pair, $\pi$ a representation of $G_f(I')$, $y$ a vector of the level and
level character of~$O$, and \m{k\in K_0(I')} with reductions $\bar k_w$ at the places $w$ of~$O$. Suppose that
at each chosen place $w$ of $O$ the columns of $\bar k_w$ indexed by the image of the embedding of~$O$ vanish
outside the rows of that image, where they form a matrix~$s_w$, and that at each twist place the same holds for
the transpose. Then \m{\pi(k)\,y=\prod_w\psi_{O,w}(\det s_w)^{-1}\cdot\rho_O\bigl((s_w)_{w\ \mathrm{chosen}}\bigr)\,y},
with $\psi_{O}$ the tame part of the character of~$O$ and $\rho_O$ as in Lemma~\ref{l:type}.}

\vspace{3pt}\phantomsection\label{i:8}\textbf{Input 8} (the data; Chapter III and the unit on the exchange).\enspace
\emph{\textup{(a)} For every infinite place~$\tau$, the vector $l_D(\tau)$ has gap $h(D)$ for the signs
\m{\fs^{Y}_{D,0}}. \textup{(b)} If \m{m(D)\le i}, all signs \m{\fs^{Y}_{D,i}} are~$1$; if \m{i<m(D)}, exactly
\m{\lfloor(m(D)-i)/2\rfloor} of them are~$-1$. \textup{(c)} \m{500\,r(D)^3+10\,r(D)^5\le h(D)}, \m{500\,r(D)^3\le h_2(D)}
and \m{16\,[E(D):\mathbb Q]\,(r(D)+1)^3\le h_2(D)}. \textup{(d)} For a chosen prime $v$ and an
integer $n$ with \m{0<|n|\le8m(D)}: \m{\omega_{D,v}^{\,n}\ne1}. \textup{(e)} The set of the primes dividing
$\pi_N(D)$ contains no conjugate of a member, and $\pi_N(D)$ has valuation~$1$ at each of them.}}
\leanline{\lean{commuteprep}\sep \lean{iSupeigenPart}\sep \lean{finiteeigenPart}\sep \lean{existsheckeproj}\sep
\lean{occurringLabelssubsingleton}\sep \lean{auxLift}\sep \lean{auxLiftinjective}\sep \lean{liftsauxLift}\sep
\lean{auxLiftprincipalRep}\sep \lean{iotaAtLevel}\sep \lean{iotaAtLevelinjective}\sep
\lean{existsisSphericalWithiota}\sep \lean{algFormRepiota}\sep \lean{smulchosenGLRepeq}\sep
\lean{hasGapgridParam}\sep \lean{cardneggridSigns}\sep \lean{omegazpowneone}}

Input~6 is the theorem of the note \emph{At most one primitive label for a Hecke eigensystem}, where it is
proved from two results on Galois representations. Further, the theorem on the cubic twist of the note \emph{Primitive labels}, the
cubic twist, is used in §\ref{s:hecke}; it belongs to this unit.

\sect{The spaces and their forms}{s:spaces}
\begin{proof}[Proof of Lemma~\ref{l:Y}]
If \m{m(D)\le i}, the signs \m{\fs^{Y}_{D,i}} are all~$1$ by Input~8(b), and the diagonal space of these signs is
standard. Let \m{i<m(D)}, and put \m{\beta=\indexY{D}{i}{i<m(D)}}, \m{J=\mathrm{Gr}(\beta,N(D))} and
\m{s=2\varrho(D)+1-i=m(D)-1-i}, the length of the string. By Input~8(b) the signs have
\m{\lfloor(s+1)/2\rfloor} entries~$-1$, so by Input~3(a) there is, for every infinite place~$\tau$, a complex
matrix $y_\tau$ with \m{y_\tau^*\,\tau(J)\,y_\tau=\mathrm{diag}(\fs^{Y}_{D,i})}; it is invertible. Let
\m{B=\mathrm{Lb}(\beta,\ideleA{D},\piIdele{D},\latticeScale{D}{i})}. By Lemma~\ref{l:pi} and Input~3(b) the matrix
\m{B^*JB} is hermitian with entries and inverse over~$\hat O_E$, so by Input~3(c) there is
\m{k_1\in\GL_{I_i(D)}(\hat O_E)} with \m{(Bk_1)^*J(Bk_1)=1}. Let \m{k_2\in K_0(I_i(D))} be the library's lift of the
family of the inverses of the reductions of $k_1$ at the chosen primes, and \m{k=k_1k_2}. Then $k$ has
reduction~$1$ at every chosen prime, and \m{(Bk)^*J(Bk)=k_2^*k_2=1}. The pair of $J$ and the adelic matrix with
components $y_\tau$ and finite part $Bk$ is the space required.
\end{proof}

\begin{lemma}[the weights are far from the row parameter]\label{l:gap}
In Setting~\ref{s:data}, let $i$ and $n$ be natural numbers and $\tau$ an infinite place of~$E(D)$. Then
\begin{align*}
\varkappa_D(\tau,n)+w_{D,i}(\tau)&\ge h_2(D)&&\text{if }n<k(D)+i-1,\\
\varkappa_D(\tau,n)+w_{D,i}(\tau)&\le-h_2(D)&&\text{otherwise.}
\end{align*}
\end{lemma}
\begin{proof}
If \m{n+1<k(D)}, then \m{n<k(D)+i-1}, and, since \m{k(D)\le r(D)\le r(D)^8},
\[\varkappa_D(\tau,n)=\lfloor k(D)/2\rfloor-1-n\ \ge\ -k(D)\ \ge\ -h_2(D);\]
as \m{w_{D,i}(\tau)\ge h(D)=2h_2(D)}, the sum is at least~$h_2(D)$.

Let now \m{n+1\ge k(D)}, and put \m{i'=n+2-k(D)}, so that \m{i'\ge1}. Then
\[\varkappa_D(\tau,n)+w_{D,i}(\tau)=h_2(D)+w_{D,i}(\tau)-w_{D,i'}(\tau)=h_2(D)+(i-i')\,h(D).\]
The condition \m{n<k(D)+i-1} means \m{i'\le i}, and then the sum is at least $h_2(D)$; otherwise \m{i'\ge i+1} and
it is at most \m{h_2(D)-h(D)=-h_2(D)}.
\end{proof}
\leanline{\lean{varsigmaaddwgap}}

\begin{lemma}[gaps]\label{l:hasgap}
In Setting~\ref{s:data}, let $\tau$ be an infinite place of~$E(D)$. Then the vector $l_D(\tau)$ has gap $h(D)$
for the signs \m{\fs^{Y}_{D,0}} of \m{\spaceY{D}{0}}; and for every natural number~$i$ it has gap $h(D)$ for the
signs \m{\fs^{W}_{D,i}} of \m{\spaceW{D}{i}}.
\end{lemma}
\begin{proof}
The first assertion is Input~8(a). That the entries of $l_D(\tau)$ are distinct and of the right parity does
not involve the signs. Let \m{\iota,\iota'\in I(D)} have different signs in \m{\fs^{W}_{D,i}}, say
\m{\num_D(\iota)<k(D)+i-1\le\num_D(\iota')}. By Lemma~\ref{l:gap} the two values of
\m{\varkappa_D(\tau,\cdot)+w_{D,i}(\tau)} are \m{\ge h_2(D)} and \m{\le-h_2(D)}, so the two entries of $l_D(\tau)$
differ by at least \m{4h_2(D)=2h(D)}.
\end{proof}
\leanline{\lean{hasGapsignsYparamW}\sep \lean{hasGapparamW}}

\begin{lemma}[cusp kernel, finite dimension]\label{l:cusp}
In Setting~\ref{s:data}, let $\fs$ be a family of signs on $I(D)$ such that for every infinite place~$\tau$ the
vector $l_D(\tau)$ has gap $h(D)$ for~$\fs$. Then \m{\epsilon_{D,\fs}} is a cusp kernel of order~$h_2(D)$, and for
every space $\cV$ over $E(D)$ with index set $I(D)$ and signs~$\fs$ the space \m{\Prin{D}{\cV}} has finite
dimension.
\end{lemma}
\begin{proof}
By Input~8(c), \m{500r(D)^3+h_2(D)\le2h_2(D)=h(D)}, and a vector with gap $h(D)$ has every smaller gap; so
Input~2(a) applies with \m{B=h_2(D)}. The group $K_g(D)$ is open, so Input~2(b) gives the finite dimension.
\end{proof}
\leanline{\lean{isCuspKernelparamW}\sep the finite dimension is an instance without a name}

By Lemmas \ref{l:hasgap} and~\ref{l:cusp}, the spaces \m{\spaceY{D}{0}} and \m{\spaceW{D}{i}} satisfy the
hypothesis on gaps of Setting~\ref{s:two}, and their spaces of principal forms have finite dimension.

\sect{Levels, moduli and auxiliary pairs}{s:levels}
\begin{lemma}\label{l:Q}
In Setting~\ref{s:data}, let $\vartheta$ be a tame character of~$D$ and $v$ a prime of~$E(D)$. If $v$ or
$\bar v$ belongs to $Q(\vartheta)$, then the product of the norms of the members of $Q(\vartheta)$ lies in~$v$,
and neither $\overline{\pi_g(D)}$ nor $N(D)$ lies in~$v$.
\end{lemma}
\begin{proof}
The norm of a prime lies in it and in its conjugate, which gives the first assertion. So the product of the
norms lies in~$v$, and it is prime to \m{g_0(D)g_1(D)N(D)}. Hence \m{g_0(D)g_1(D)N(D)} does not lie in~$v$: otherwise
\m{v\cap\mathbb Z} would contain two coprime integers, hence~$1$. As
\m{\overline{\pi_g(D)}\pi_g(D)=g_0(D)g_1(D)}, neither $\overline{\pi_g(D)}$ nor $N(D)$ lies in~$v$.
\end{proof}
\leanline{\lean{levelNormmemandnotMemofmemQ}}

\begin{definition}[places and moduli of a tame character]\label{d:levels}
In Setting~\ref{s:data}, let $\vartheta$ be a tame character of~$D$.
\begin{enumerate}[label=\textup{(\alph*)},leftmargin=2.3em,itemsep=2pt]
\item \m{\levelPlaces{D}{\vartheta}} is the system of places with the chosen primes $\Ch(D)$ as its chosen places
and the members of $Q(\vartheta)$ as its twist places. (It is a system of places in the library's sense: these
primes are distinct, and none is the conjugate of one of them, by Lemma~\ref{l:Q}.)
\item For a natural number $n$ and a Boolean~$o$, \m{\levelMod{D}{\vartheta}{n}{o}} is the ideal of $O_{E(D)}$
\[
\Bigl(\prod_{v\in\Ch(D)}v\Bigr)^{n}\cdot\prod_{v\in Q(\vartheta)}v\quad\text{if $o$ is true,}\qquad
\Bigl(\prod_{v\in\Ch(D)}v\Bigr)^{n}\quad\text{if $o$ is false.}
\] The formal text takes $o$ true when it is not given.
\end{enumerate}
\end{definition}
\leanline{\lean{levelPlaces}\sep \lean{decl}}

\begin{lemma}\label{l:nonempty}
In Setting~\ref{s:data}, let $\vartheta$ be a tame character of~$D$, $n$ a natural number with \m{n\ge1}, and $o$
a Boolean. Then some prime of $E(D)$ divides \m{\levelMod{D}{\vartheta}{n}{o}}.
\end{lemma}
\begin{proof}
There is a chosen prime $v$ above~$g_0(D)$. The ideal is contained in \m{(\prod_{v'\in\Ch(D)}v')^n}, which, as
\m{n\ge1}, is contained in the product of the chosen primes and so in~$v$.
\end{proof}
\leanline{\lean{nonemptymodulusplaces}}

\begin{definition}[the auxiliary pairs]\label{d:aux}
An \emph{auxiliary pair} over a CM field, in the sense of the unit on the rows, is a record of: a system of
places; a space with an index set $I_1$ and signs; an embedding of $I_1$ into a larger index set~$I'$; a numbering
\m{I_1\to\mathbb N}; two families of natural numbers \m{\kappa^{+}(\tau,n)}, \m{\kappa^{-}(\tau,n)} (numbers of columns of
height~$n$, which give the weight of forms on the standard space with index set~$I'$); and a character of the
norm-one ideles with a weight and with a tame part that is trivial off the places of the system.

In Setting~\ref{s:data}, let $\vartheta$ be a tame character of~$D$ and $i$ a natural number. The auxiliary pair
\m{\auxPair{D}{\vartheta}{i}} has: the system of places \m{\levelPlaces{D}{\vartheta}}; the space
\m{\spaceW{D}{i}} with its signs \m{\fs^{W}_{D,i}}; the inclusion of $I(D)$ into
\m{I(D)\sqcup\{0,\dots,t(D)-1\}}; the numbering \m{\iota\mapsto\num_D(\iota)} where the sign is~$1$ and
\m{\iota\mapsto r(D)-1-\num_D(\iota)} where it is~$-1$; the column numbers, written
\m{\kappa^{+}_{O}(\tau,n)} and \m{\kappa^{-}_{O}(\tau,n)} with \m{O=\auxPair{D}{\vartheta}{i}} as in the note on the
auxiliary pair,
\begin{align*}
\kappa^{+}_{O}(\tau,n)&=\Bigl(\varkappa_D(\tau,n-1)-\max\bigl(\varkappa_D(\tau,n),\lfloor t(D)/2\rfloor-w_{D,i}(\tau)\bigr)-1\Bigr)_+,\\
\kappa^{-}_{O}(\tau,n)&=\Bigl(\min\bigl(\varkappa_D(\tau,r(D)-n-1),-w_{D,i}(\tau)-\lfloor t(D)/2\rfloor\bigr)\\
&\qquad\qquad-\varkappa_D(\tau,r(D)-n)-1\Bigr)_+,
\end{align*}
where \m{(\cdot)_+} is the positive part and the differences \m{n-1}, \m{r(D)-n-1}, \m{r(D)-n} of natural numbers are
truncated at zero; and the character \m{\phi_{D,i}\,\vartheta}, where $\phi_{D,i}$ is the character of the data of
weight \m{\tau\mapsto-w_{D,i}(\tau)} with tame part \m{(\omega_{D,v}^{\,i})_v}, so that the tame part of the
product is \m{(\omega_{D,v}^{\,i}\,\psi_{\vartheta,v})_v}.
\end{definition}
\leanline{\lean{auxPair}}

\sect{The good idempotent}{s:idem}
\begin{lemma}\label{l:central}
In Setting~\ref{s:data}, the element $e_0(D)$ is an idempotent of \m{\mathbb C[\sG(D)]} whose coefficients are
constant on conjugacy classes.
\end{lemma}
\begin{proof}
It is a sum of elements \m{e^{\otimes}_{I(D),\mathbf c}} over a finite set of families; apply Input~4(a).
\end{proof}
\leanline{\lean{isCentralProje}}

So $e_0(D)$ is central, and \m{\prep{D}{\cV}(e_0(D))} is a projector of \m{\Prin{D}{\cV}} onto \m{\Good{D}{\cV}}
that commutes with \m{\prep{D}{\cV}}.

\begin{lemma}[good forms see no reflection characters]\label{l:refl}
In Setting~\ref{s:data}, let \m{(y_v)_{v\in\Ch(D)}} be a family of elements \m{y_v\in\mathbb C[\GL_{I(D)}(\F_v)]},
let $v_0$ be a chosen prime, $\chi$ a multiplicative character of~$\F_{v_0}$ and \m{l\in\F_{v_0}^{I(D)}} a vector.
Suppose that \m{\chi^{M_\omega(D)}=1}, that \m{l\ne0}, and that
\m{y_{v_0}=\sum_h\chi(\det h)^{-1}h}, the sum over the \m{h\in\GL_{I(D)}(\F_{v_0})} of the form
\m{1+u\,l^{\mathsf T}}. Then \m{\bigl(\bigotimes_{v}y_v\bigr)\cdot e_0(D)=0}.
\end{lemma}
\begin{proof}
It suffices to show \m{(\bigotimes_vy_v)\cdot e^{\otimes}_{I(D),\mathbf c}=0} for every family $\mathbf c$ in the
sum that defines~$e_0(D)$. This product is the tensor product of the \m{y_v\,e_{I(D),c_v}}. By
Definition~\ref{d:good}, no \m{x\in\F_{v_0}} with \m{x^{M_\omega(D)}=1} is a root of~$c_{v_0}$; and
\m{x_\chi^{M_\omega(D)}=1}, because \m{\chi\mapsto x_\chi} is a homomorphism. So $x_\chi$ is not a root of
$c_{v_0}$, and the factor at $v_0$ vanishes by Input~4(b).
\end{proof}
\leanline{\lean{reflectionSummule}}

\sect{Hecke operators: one operator for two pairings; equal traces}{s:hecke}
For a tame character $\vartheta$ of $D$, a natural number~$i$ and a space $\cV$ over $E(D)$ with index set
$I(D)$, the unit on the exchange defines the \emph{Rallis terms}. Call a prime of $E(D)$ \emph{bad} if it or its
conjugate divides \m{\levelMod{D}{\vartheta}{1}{\mathrm{true}}}; the chosen primes are bad. Let
\m{L_{D,\vartheta}} be the subgroup of the \m{g\in G_f(I(D))} whose components at the primes that are not bad
are integral. The Rallis terms are the numbers, for \m{z,y\in\Prin{D}{\cV}} and
\m{\gamma\in G_f(I(D))/L_{D,\vartheta}},
\[
\mathrm{Ral}_{D,\vartheta,i}(\cV;z,y;\gamma)=c_{D,\vartheta,i}(\gamma)\,\langle z,R(g_\gamma)\,y\rangle .
\]
Here \m{g_\gamma\in G_f(I(D))} is a representative of $\gamma$, chosen by the library, whose components at the
bad primes are~$1$, and \m{c_{D,\vartheta,i}(\gamma)} is a complex number
that does not depend on the space: a matrix coefficient of the vacuum vector of a Weil representation at
$g_\gamma$, for the character of \m{\auxPair{D}{\vartheta}{i}}, divided by a power of the norm of that modulus.

\begin{lemma}[one Hecke operator for two pairings]\label{l:common}
In Setting~\ref{s:data}, let $\fs$, $\fs'$, $\cV$, $\cV'$ be as in Setting~\ref{s:two}. Let $\vartheta$ be a tame
character of~$D$, $i$ a natural number, $A$ and $A'$ complex numbers, $B$ a complex function on
\m{\Prin{D}{\cV}\times\Prin{D}{\cV}} and $B'$ one on \m{\Prin{D}{\cV'}\times\Prin{D}{\cV'}}. Suppose that for all
\m{z,y\in\Prin{D}{\cV}} the family \m{\bigl(A\cdot\mathrm{Ral}_{D,\vartheta,i}(\cV;z,y;\gamma)\bigr)_\gamma} is
summable with sum \m{B(z,y)}, and that for all \m{z,y\in\Prin{D}{\cV'}} the family
\m{\bigl(A'\cdot\mathrm{Ral}_{D,\vartheta,i}(\cV';z,y;\gamma)\bigr)_\gamma} is summable with sum \m{B'(z,y)}. Then
there is \m{T\in\cH(D)} such that
\begin{align*}
B(z,y)&=A\,\langle z,T_\cV\,y\rangle&&\text{for all }z,y\in\Prin{D}{\cV},\\
B'(z,y)&=A'\,\langle z,T_{\cV'}\,y\rangle&&\text{for all }z,y\in\Prin{D}{\cV'}.
\end{align*}
\end{lemma}
\begin{proof}
For \m{T\in\cH(D)} let $\Psi(T)$ be the pair of the function \m{(z,y)\mapsto A\langle z,T_\cV y\rangle}
on pairs of forms on~$\cV$ and of the function \m{(z,y)\mapsto A'\langle z,T_{\cV'}y\rangle} on pairs of forms
on~$\cV'$. The map $\Psi$ is linear, and it factors through the pair
\m{(T_\cV,T_{\cV'})} of endomorphisms of two spaces of finite dimension (Lemma~\ref{l:cusp}); so its image
$\mathcal Z$ is a subspace of finite dimension of the space of pairs of functions, hence closed for pointwise
convergence. The pair \m{(B,B')} is the limit of the finite partial sums of the pairs of terms. So it suffices
to show that for every coset~$\gamma$ the pair of the two terms lies in~$\mathcal Z$.

Fix $\gamma$, write \m{g=g_\gamma}, and let $H$ be the subgroup of the \m{k\in K_g(D)} with
\m{g^{-1}kg\in K_g(D)}. It is open in the compact group~$K_g(D)$, so of finite index. Put
\[
T^\gamma=\frac{c_{D,\vartheta,i}(\gamma)}{[K_g(D):H]}\sum_{k\in K_g(D)/H}R(kgk^{-1}).
\]
The elements \m{kgk^{-1}} have component~$1$ at every chosen prime, because $g$ has: the chosen primes
are bad. For a form $y$ fixed by~$K_g(D)$, \m{R(kgk^{-1})y=R(k)R(g)y}; this depends only on the class
of $k$ modulo~$H$, and the sum over the classes is again fixed by~$K_g(D)$. So $T^\gamma$ maps
\m{\Prin{D}{\cV}} to itself for every space~$\cV$, and \m{T^\gamma\in\cH(D)}. As $R$ is unitary and $z$ is fixed
by~$K_g(D)$,
\[
\langle z,T^\gamma_\cV\,y\rangle=\frac{c_{D,\vartheta,i}(\gamma)}{[K_g(D):H]}\sum_k\langle R(k)^{-1}z,R(g)y\rangle
=\mathrm{Ral}_{D,\vartheta,i}(\cV;z,y;\gamma),
\]
and the same holds for~$\cV'$ with the same~$T^\gamma$. So the pair of terms is \m{\Psi(T^\gamma)}.
\end{proof}
\leanline{\lean{existscommonheckerallis}}

\begin{proof}[Proof of Theorem 1]
By Lemma~\ref{l:cusp}, \m{\Prin{D}{\cV}} and \m{\Prin{D}{\cV'}} have finite dimension. Let $\mathcal F$ be the
set of the eigensystems $p$ with \m{\Prin{D}{\cV}[p]\ne0} or \m{\Prin{D}{\cV'}[p]\ne0}; it is finite, and each of
the two spaces is the sum of its parts \m{[p]}, \m{p\in\mathcal F} (Input~5(b)). For \m{p\in\mathcal F} let
\m{\Pi^p\in\cH(D)} be as in Input~5(c), for $p$ and~$\mathcal F$.

\emph{Twisting the labels.} Fix \m{p\in\mathcal F} and a chosen prime~$v$. The sets \m{\Occ(D,\cV,p,v)} and
\m{\Occ(D,\cV',p,v)} have at most one element each, by the two hypotheses of rigidity; $I(D)$ has $r(D)$ elements
and $r(D)$ is a power of~$3$ greater than~$1$; and $9$ divides \m{\#\F_v-1}. So the theorem on the cubic twist of the note
\emph{Primitive labels} gives, for every polynomial \m{c\in\F_v[X]}, a polynomial \m{c^{p,v}\in\F_v[X]} such that:
if $c$ is a primitive label of degree~$r(D)$, then so is $c^{p,v}$, and $c^{p,v}$ lies in neither of the two sets;
and \m{e_{I(D),c}(\gamma)=e_{I(D),c^{p,v}}(\gamma)} for every \m{\gamma\in\GL_{I(D)}(\F_v)} whose characteristic
polynomial is not irreducible. Put
\[
\tilde e_p=\sum_{\mathbf c}e^{\otimes}_{I(D),(c_v^{p,v})_v}\ \in\mathbb C[\sG(D)],
\]
the sum over the same families \m{\mathbf c=(c_v)_v} as in the definition of~$e_0(D)$.

\emph{The twisted idempotent kills the part $[p]$.} Let $\cV''$ be $\cV$ or~$\cV'$, and $\mathbf c$ a family in
that sum. Above $g_0(D)$ there is a chosen prime $v$ at which $c_v$ is a primitive label. Then $c_v^{p,v}$ is a
primitive label of degree~$r(D)$ outside \m{\Occ(D,\cV'',p,v)}, which by the definition of that set means that
\m{e_{I(D),c_v^{p,v}}}, acting through the factor~$v$, kills \m{\Prin{D}{\cV''}[p]}. By Input~4(c) the tensor
product kills it too. Hence \m{\prep{D}{\cV''}(\tilde e_p)} is zero on \m{\Prin{D}{\cV''}[p]}.

\emph{A formula for the trace.} On \m{\Prin{D}{\cV''}},
\begin{equation}\label{e:sum}
T_{\cV''}\circ\prep{D}{\cV''}\bigl(e_0(D)\bigr)=\sum_{p\in\mathcal F}(T\,\Pi^p)_{\cV''}\circ\prep{D}{\cV''}\bigl(e_0(D)-\tilde e_p\bigr).
\end{equation}
It suffices to check this on each part \m{\Prin{D}{\cV''}[p']}, \m{p'\in\mathcal F}. The representation
\m{\prep{D}{\cV''}} preserves that part (Input~5(a)); on it $\Pi^p$ is the identity for \m{p=p'} and zero for the
other \m{p\in\mathcal F}; and \m{\prep{D}{\cV''}(\tilde e_{p'})} vanishes there. So both sides of \eqref{e:sum}
are \m{T_{\cV''}\circ\prep{D}{\cV''}(e_0(D))} on that part. Now \m{\prep{D}{\cV''}(e_0(D))} is a projector onto
\m{\Good{D}{\cV''}} that commutes with $T_{\cV''}$ (Lemma~\ref{l:central}, Input~5(a)), so the trace of the left
side of \eqref{e:sum} is the trace of $T$ on \m{\Good{D}{\cV''}}.

For \m{p\in\mathcal F} write \m{(T\,\Pi^p)_{\cV''}=\sum_gc^p_g\,R(g)}, a finite sum over elements
\m{g\in G_f(I(D))} with component $1$ at the chosen primes, the same for $\cV$ and~$\cV'$. Let \m{(b_n)_n} be an
orthonormal basis of \m{\Prin{D}{\cV''}}. The trace of an endomorphism $S$ is \m{\sum_n\langle b_n,Sb_n\rangle},
so by \eqref{e:sum} and the definition of \m{\prep{D}{\cV''}}
\begin{equation}\label{e:trace}
\mathrm{tr}\bigl(T\mid\Good{D}{\cV''}\bigr)=\sum_{p\in\mathcal F}\ \sum_{\gamma\in\sG(D)}\ \sum_g\,(e_0(D)-\tilde e_p)(\gamma)\;c^p_g\;
\sum_n\langle b_n,R(g\tilde\gamma)\,b_n\rangle .
\end{equation}

\emph{Only elliptic elements remain.} Let \m{\gamma=(\gamma_v)_v\in\sG(D)} be such that no $\gamma_v$ has an
irreducible characteristic polynomial. Then for every family $\mathbf c$ and every~$v$,
\m{e_{I(D),c_v}(\gamma_v)=e_{I(D),c_v^{p,v}}(\gamma_v)}; taking products over $v$ and summing over $\mathbf c$ gives
\m{(e_0(D)-\tilde e_p)(\gamma)=0}. So in \eqref{e:trace} only those $\gamma$ contribute for which some
$\gamma_{v_0}$ has an irreducible characteristic polynomial. For such a $\gamma$, for $g$ with component~$1$ at
the chosen primes and for \m{k\in K_g(D)}, the component at $v_0$ of \m{g\tilde\gamma k} is that of
\m{\tilde\gamma k}: a matrix over $O_{v_0}$ with reduction~$\gamma_{v_0}$. Its characteristic polynomial is monic
over~$O_{v_0}$ with irreducible reduction, hence irreducible over $O_{v_0}$ and, by Gauss's lemma, over~$E_{v_0}$.

\emph{Comparison.} So Input~1 applies to \m{\mathbf K=K_g(D)} and the element \m{g\tilde\gamma}, with \m{l=l_D},
\m{B=h_2(D)} and \m{H=h(D)}: the field $E(D)$ contains a primitive root of unity of order~$5040$ and has an even
number of infinite places, and the two inequalities hold by Input~8(c). For both $\cV$ and $\cV'$ the type is a
cusp kernel of order $h_2(D)$ (Lemma~\ref{l:cusp}) and $l_D(\tau)$ has gap~$h(D)$ (Setting~\ref{s:two}). Hence the
inner sum \m{\sum_n\langle b_n,R(g\tilde\gamma)b_n\rangle} of \eqref{e:trace} has the same value for $\cV$ and
for~$\cV'$. All other factors of \eqref{e:trace} do not depend on the space, and the theorem follows.
\end{proof}

\begin{proof}[Proof of Theorem 2]
(a) For \m{n\ge0}, \m{T^n\in\cH(D)}, and the action on good forms respects products. By Theorem~1 the $n$-th
powers of the two restrictions of $T$ have equal traces, for all \m{n\ge0}. By Input~4(d) they have the same
characteristic polynomial, hence the same determinant.

(b) Let $\mathcal F'$ be the finite set of the eigensystems $p'$ with \m{\Prin{D}{\cV'}[p']\ne0}, and $\Pi$ as in
Input~5(c) for $p$ and~$\mathcal F'$. On \m{\Prin{D}{\cV'}}, which is the sum of its parts \m{[p']},
\m{p'\in\mathcal F'}, the image of $\Pi_{\cV'}$ lies in \m{\Prin{D}{\cV'}[p]}. Since $\Pi_{\cV'}$ preserves the
good forms, it maps \m{\Good{D}{\cV'}} into \m{\Good{D}{\cV'}\cap\Prin{D}{\cV'}[p]=0}. So \m{1-\Pi} is the identity
on \m{\Good{D}{\cV'}}, with determinant~$1$; by (a) its determinant on \m{\Good{D}{\cV}} is~$1$, and
\m{1-\Pi_\cV} is injective there. But it kills \m{\Good{D}{\cV}\cap\Prin{D}{\cV}[p]}, on which $\Pi_\cV$ is the
identity. So this intersection is zero.

For the last assertion apply Theorem~1 to \m{\cV=\spaceY{D}{0}} and \m{\cV'=\spaceW{D}{i}}, which satisfy
Setting~\ref{s:two} by Lemma~\ref{l:hasgap} and have rigid labels by Theorem~3, and to \m{T=1}: the trace of the
identity is the dimension.
\end{proof}

\sect{Rigid labels: proof of Theorem 3}{s:rigid}
\begin{lemma}[the parahoric type of a vector of the level of an auxiliary pair]\label{l:type}
In Setting~\ref{s:data}, let $O$ be an auxiliary pair over $E(D)$ whose space has the index set $I_1$ and whose
embedding is \m{\nu_O\colon I_1\to I'}, with $I_1$ and $I'$ finite sets. Its \emph{level group} $K_O$ consists of
the \m{k\in K_0(I')} whose reduction fixes the basis vectors indexed by \m{\nu_O(I_1)} at the chosen places of~$O$
and preserves the span of the other basis vectors at its twist places; its \emph{level character} is
\m{k\mapsto\prod_w\psi_{O,w}^{-1}(\det s_w(k))}, the product over the twist places, with $s_w(k)$ the block of the
reduction of $k$ at $w$ indexed by \m{\nu_O(I_1)} and \m{(\psi_{O,w})_w} the tame part of the character of~$O$. Let
$\mathcal S$ be a complex vector space of functions on $G_f(I')$ with values in the functions on a finite set,
and $\pi$ a representation of $G_f(I')$ on~$\mathcal S$. On the subspace of the \m{u\in\mathcal S} on which $K_O$
acts through the level character, the group \m{\prod_w\GL_{I_1}(\F_w)}, $w$ the chosen places of~$O$, acts by
\m{\rho_O(\gamma)=\prod_w\psi_{O,w}(\det\gamma_w)\cdot\pi(\tilde\gamma)}, with $\tilde\gamma$ a lift of the block
diagonal matrices with blocks $\gamma_w$ and~$1$. Let $u$ be a vector of that subspace, $v$ a chosen place of~$O$,
which is split, and $f$ a function on $G_f(I')$ with values in the functions on that finite set. Suppose that
\[
\pi(k^{(v)})\,u=R(k^{(v)})\,f\qquad\text{for every }k\in\GL_{I'}(E_v).
\]
Then for every \m{x\in\GL_{I_1}(\F_v)} and every \m{k\in\GL_{I'}(O_v)} whose reduction $\bar k$ is block upper
triangular for the two blocks \m{\nu_O(I_1)} and its complement, with first block equal to~$x$:
\[
R(k^{(v)})\,f=\psi_{O,v}(\det x)^{-1}\cdot\rho_O(x_{(v)})\,u ,
\]
where $x_{(v)}$ has component $x$ at $v$ and $1$ at the other chosen places of~$O$. In the language of Input~6:
$f$ satisfies its condition \textup{(iii)} for the map $\beta$ that is $0$ on \m{\nu_O(I_1)} and $1$ elsewhere,
\m{n_1=1}, \m{\chi_1=1}, \m{\psi=\psi_{O,v}^{-1}}, and \m{x\mapsto\rho_O(x_{(v)})\,u} in the place of
\m{x\mapsto L(\prep{D}{\cV}(x_{(v)})y)}.
\end{lemma}
\begin{proof}
Let $x$ and $k$ be as stated. By hypothesis \m{R(k^{(v)})f=\pi(k^{(v)})u}. The reductions of $k^{(v)}$ at the
places of $O$ are $\bar k$ at $v$ and~$1$ elsewhere. Since $\bar k$ is block upper triangular with first
block~$x$, its columns indexed by \m{\nu_O(I_1)} vanish outside the rows of \m{\nu_O(I_1)} and form the
matrix~$x$ there; at the other places the condition of Input~7(c) holds with the identity matrix. So Input~7(c)
gives \m{\pi(k^{(v)})u=\psi_{O,v}(\det x)^{-1}\rho_O(x_{(v)})u}.
\end{proof}
\leanline{\lean{hasParahoricTypechosenGLRep}}

\begin{proof}[Proof of Theorem 3]
Let $p$ be an eigensystem and $v$ a chosen prime. In both cases Input~6 is applied; the gaps that it does not
ask for are not needed.

\emph{The space \m{\spaceW{D}{i}}, \m{1\le i\le m(D)}.} Take \m{I'=I(D)\sqcup\{0,\dots,t(D)-1\}}, $\nu$ the inclusion,
and \m{L=\Lambda_i} of Input~7(a); (i) and (ii) are part of that input. For (iii), Lemma~\ref{l:type} is applied
to \m{O=\auxPair{D}{1}{i}}, to the representation of right translations on the functions on $G_f(I')$, to
\m{u=f=\Lambda_i(y)}, for which its hypothesis is trivially true, and to~$v$: with the equivariance of
$\Lambda_i$ in Input~7(a) it gives (iii) for \m{n_1=1}, \m{\chi_1=1} and \m{\psi=\psi_{O,v}^{-1}}. The tame part
of the character of this pair at $v$ is \m{\omega_{D,v}^{\,i}}, because the trivial tame character has trivial
tame part; so \m{\psi=\omega_{D,v}^{-i}} and \m{\psi^{M_\omega(D)}=1}. A single character is a family of distinct
characters, and \m{1^{M_\omega(D)}=1}; this is (iv).

\emph{The space \m{\spaceY{D}{0}}.} Take \m{I'=I_{m(D)}(D)}, $\nu$ the inclusion of $I(D)$, and \m{L=\Lambda_{\mathrm{ch}}} of
Input~7(b), which gives (i) and (ii). Let $\beta$ be the map that is $0$ on $I(D)$ and $n$ on
\m{\{n\}\times\{0,\dots,t(D)-1\}}, \m{n_1=m(D)}, \m{\chi_n=\omega_{D,v}^{\,n}} and \m{\psi=1}. For
\m{k\in\GL_{I'}(O_v)} with reduction block upper triangular for $\beta$ and block~$0$ equal to~$x$, the element
$k^{(v)}$ lies in $K_0(I')$ and has reductions $\bar k$ at $v$ and $1$ at the other chosen primes; so the formula
of Input~7(b), in which all factors at the other chosen primes are~$1$, is (iii). The characters
\m{\omega_{D,v}^{\,n}}, \m{1\le n\le m(D)}, are distinct: a quotient of two of them is \m{\omega_{D,v}^{\,n-n'}} with
\m{0<|n-n'|\le8m(D)}, which is not trivial by Input~8(d). They have order dividing \m{M_\omega(D)}, the order
of~$\omega_{D,v}$. This is (iv).
\end{proof}

\sect{Where the results are used}{s:use}
The definitions of this note are used throughout Chapters V to~VIII: \m{\Row{D}} is the space of rows of the
three main theorems, and \m{\spaceY{D}{i}}, \m{\spaceW{D}{i}} are the spaces of the chain of theta lifts and of the
single lifts. Theorem~2 enters the proof of Theorem~C in three places. Its part~(a), with Lemma~\ref{l:common},
is the step that writes the volume of the chain of lifts on the rows as a constant times the product of the
volumes of the single lifts (unit on the rows): by Lemma~\ref{l:common} the two Rallis inner products are given
by one Hecke operator, and by part~(a) its determinants on \m{\Row{D}} and on \m{\Good{D}{\spaceW{D}{i}}} agree.
The equality of dimensions is used there too, and in Chapter~VI to bound \m{\rowDim{D}} from below. Part~(b)
is used in the construction of the fit family of twists (Chapter~VIII), to pass from eigensystems of good forms
on~\m{\spaceW{D}{i}} to those of rows. Lemma~\ref{l:refl} is used in the comparison of forms on the rows, and
Lemma~\ref{l:type} again in Chapter~VIII. Lemmas \ref{l:m}, \ref{l:pi}, \ref{l:Q} and~\ref{l:nonempty} are used in
the units on the seesaw and on the exchange.

\sect{Notes on the formal text}{s:notes}
\begin{enumerate}[label=\textup{\arabic*.},leftmargin=1.6em,itemsep=5pt]
\item \emph{Letters.} The formal text writes {\leanfont V}, {\leanfont s} for a space and its signs,
{\leanfont j} for the index~$i$, {\leanfont w} and {\leanfont x} for primes and chosen primes, {\leanfont p} for a
Hecke element in the file on traces and for an eigensystem in the file on rigidity, {\leanfont q} and
{\leanfont d} for eigensystems, {\leanfont a}, {\leanfont a'} for the constants $A$, $A'$, {\leanfont c},
{\leanfont b} for the exponent $n$ and the Boolean $o$ of a modulus, {\leanfont ς} for the row parameter and
{\leanfont paramW} for its double~$l_D$. The spaces \m{\spaceY{D}{i}}, \m{\spaceW{D}{i}} are the authors'
{\leanfont Y\_j} and {\leanfont W\_j}.
\item \emph{The hypothesis on gaps} of Setting~\ref{s:two} is an instance argument in the formal text: it is
found by Lean, not passed by hand. Its only role in Lemma~\ref{l:common} is to make the two spaces of forms
finite-dimensional.
\item \emph{Lemma~\ref{l:pi}.} The formal text states the valuations as values \m{\exp(-1)} and~$1$ of the
valuation at~$v$, written multiplicatively in the value group; this is the order $1$ and~$0$. The proof there
quotes Input~8(e).
\item \emph{Lemma~\ref{l:gap}.} The formal proof uses only the inequality \m{w_{D,i}(\tau)+h(D)\le w_{D,i'}(\tau)}
for \m{i<i'}, an input from the unit on the exchange, and closes the cases by an automatic tactic; I used the
definition of the weights, which gives the difference exactly. After the lemma the formal text registers it as a
pattern for that tactic; this line is not part of a proof.
\item \emph{Lemma~\ref{l:Y}.} In (b) the formal statement is made for every proof of \m{i<m(D)}. The formal proof
takes for the infinite components of the point the inverses of the matrices
\m{y_\tau\,\mathrm{diag}(\fs^{Y}_{D,i})}-type frames that Input~3(a) provides, arranged so that the defining
identity of a space holds; I did not trace this arrangement term by term, and the dossier of the unit marks the
same place as uncertain.
\item \emph{Theorem~1.} The formal proof first shows identity \eqref{e:sum} for an arbitrary space whose occurring
labels are among those of $\cV$ and~$\cV'$ and whose eigensystems lie in~$\mathcal F$, and then applies it twice.
It uses of a good family only that it has a primitive label above~$g_0(D)$. The Hecke elements $\Pi^p$ are chosen for
all eigensystems, and the theorem on the cubic twist of the note \emph{Primitive labels} is applied to all pairs \m{(p,v)}, not only
for \m{p\in\mathcal F}.
\item \emph{Theorem~2.} In (a) the formal proof writes the determinants as products of the roots of the
characteristic polynomials. The traces for \m{n=0} give the equality of the dimensions, which the comparison of
characteristic polynomials of matrices of different sizes needs.
\item \emph{Theorem~3.} The applications of Input~6 pass most of its arguments implicitly; in the case of
\m{\spaceY{D}{0}} two side conditions on the places other than~$v$ are closed by automatic tactics, and I read
them as the statement that the reductions of $k^{(v)}$ at those places are~$1$. The dossier marks the same lines.
\item \emph{Arguments.} For Lemma~\ref{l:Y}, Lemma~\ref{l:cusp} with its instance, and Lemma~\ref{l:refl}, the
dossier records the arguments as derived by hand from the text of the signatures, not as printed by Lean; the
second reading of this write-up asks for Lean's own output before these records are called confirmed. The
statements above agree with the derived lists: Lemma~\ref{l:Y} takes $D$ and~$i$ only; Lemma~\ref{l:cusp} takes
the hypothesis on gaps as an instance argument; Lemma~\ref{l:refl} takes the family, the character and the vector
implicitly and the chosen prime and the three hypotheses explicitly.
\item \emph{What the final theorems do not use.} The plan of the write-up, which follows names in the text of
the proofs, records most but not all of the \UnitMembers{} declarations of the unit as reached from the final
theorems. Instances are used by Lean without being named, so which declarations are really unused is for the
kernel export to say.
\end{enumerate}

\subhead{How this note was written.} From the dossier of the unit: its statements and the structured proofs of
Lemmas \ref{l:pi}, \ref{l:Y}, \ref{l:gap}, \ref{l:nonempty}, \ref{l:refl}, \ref{l:common}, \ref{l:type} and of
Theorems 1 to~3, which a program compared line by line with the source (\UnitSteps{} leaf steps in the
\UnitProofs{} proofs of the unit), and the second reading of the unit by the other agent of this write-up. I read
the four Lean files of the unit completely; in the other units the definitions of the Hecke algebra, of the
eigenparts and of the occurring labels, of the good idempotent, of the weights and of the row parameter, of
auxiliary pairs and of the Rallis terms; and in the library the definitions of spaces, forms, levels and types,
and the statement of Input~1. The statements of the other inputs I took from the dossier's quotations, which the
program compared with the source. Every statement of a declaration of this note was compared with its Lean
statement argument by argument, and the order of the quantifiers is that of the formal text. The proofs follow
the formal ones; the proof of Lemma~\ref{l:gap} is shorter, and the proof of Theorem~1 is arranged as a
computation with an explicit formula for the trace. The descriptions of the objects of Setting~\ref{s:data} are
recollections, and the notes of their units are the authority for them. The references are those of the
library's files, of the catalogue of this handout and of its second reading; I did not open the cited works.

\head{Appendix A. Concordance with the Lean artifact}
{\fontsize{8.6}{11.6}\selectfont\renewcommand{\leanfont}{\ttfamily\fontsize{7.7}{10}\selectfont\color{leancol}}
\setlength{\LTleft}{0pt}\setlength{\LTright}{\fill}\setlength{\tabcolsep}{0pt}
\begin{longtable}{@{}p{30mm}p{56mm}p{43mm}r@{}}
\textbf{In this note} & \textbf{Lean declaration} & \textbf{File} & \textbf{Lines}\\\midrule\endhead
\ConcordanceRows
\end{longtable}}

\subhead{The main statements as they are written in the artifact.}
\quoted{traceheckeOnGoodeq}
\quoted{detheckeOnGoodeq}
\quoted{finrankrowseqfinrankgoodForms}
\quoted{existscommonheckerallis}
\quoted{traceComparison}
The first three stand after a line that declares the variables {\leanfont D}, {\leanfont s}, {\leanfont s'},
{\leanfont V}, {\leanfont V'}, the two instance arguments on gaps and the two hypotheses of rigidity
{\leanfont o}, {\leanfont o'}. In these texts {\leanfont principalForms}, {\leanfont goodForms},
{\leanfont heckeAlgebra}, {\leanfont heckeOnGood}, {\leanfont rallisTerm} are \m{\Prin{D}{\cV}}, \m{\Good{D}{\cV}},
$\cH(D)$, the action of $\cH(D)$ on good forms and the Rallis terms; {\leanfont spaceW}, {\leanfont rowDim} are
\m{\spaceW{D}{i}} and \m{\rowDim{D}}; {\leanfont typeKernel l s} is \m{\prod_\tau\epsilon_{l(\tau)}},
{\leanfont formRep} is~$R$ and {\leanfont levelSpace} takes the vectors fixed by a subgroup.

\head{Notation}
\begingroup\fontsize{8.6}{10.9}\selectfont\raggedright
\setlength{\columnsep}{9mm}\setlength{\parskip}{0pt}
\begin{multicols}{2}
\nt{$P=(N_0,d)$, $D$}{the parameters; admissible data (Setting~\ref{s:data})}
\nt{$E$, $E(D)$, $K$, $O_E$}{a CM field; the field of the data; \m{\mathbb Q(\sqrt{-d})}; the ring of integers}
\nt{$\mathbf K$}{an open subgroup of $K_0(I)$ or of $G_f(I)$ (Inputs 1 and~2)}
\nt{$v$, $w$, $\bar v$, $\F_v$, $E_v$, $O_v$}{primes of a number field, the conjugate prime, residue field, completion, its integers}
\nt{$\mathbb A_E$, $\mathbb A_{E,f}$, $\hat O_E$}{adeles, finite adeles, integral finite adeles}
\nt{$\tau$}{an infinite place of $E(D)$}
\nt{$m(D)$, $r(D)$, $x(D)$, $\mu(D)$}{components of $D$: the number of lifts, the rank, the first scale, the exponent}
\nt{$k(D)$, $t(D)$, $\varrho(D)$, $n_i(D)$}{\m{r(D)+1-m(D)}, \m{5r(D)}, \m{m(D)/2-1}, \m{k(D)+i(t(D)+1)}}
\nt{$h_2(D)$, $h(D)$, $M_\omega(D)$}{\m{r(D)^8}, \m{2r(D)^8}, \m{(N_0d)^{\mu(D)+6}}}
\nt{$w_{D,i}(\tau)$, $w_D(\tau)$, $N_1(D)$}{the weights; the nominal weight of the overview; the order of the root of unity that generates $E(D)$ over $K$}
\nt{$\varkappa_D(\tau,n)$, $l_D(\tau)$}{the row parameter, its double as a vector on $I(D)$}
\nt{$N(D)$, $\pi_N(D)$, $g_0(D)$, $g_1(D)$, $\pi_g(D)$}{the level and an element of norm $N(D)$; two primes and an element of norm $g_0(D)g_1(D)$}
\nt{$\Ch(D)$, $\sG(D)$}{the chosen primes; \m{\prod_{v\in\Ch(D)}\GL_{I(D)}(\F_v)}}
\nt{$I(D)$, $\num_D$, $I_i(D)$, $\mathcal X_i(D)$}{the row index set and its numbering; the index sets of \m{\spaceY{D}{i}} and of its first summand}
\nt{$\omega_{D,v}$}{the character of order $M_\omega(D)$ of $\F_v^\times$ at a chosen prime}
\nt{$\vartheta$, $Q(\vartheta)$, $\psi_{\vartheta,v}$}{a tame character, its set of primes, its tame part}
\nt{$I$, $I'$, $I_1$, $\GL_I$, $g^*$}{finite index sets; invertible matrices indexed by $I$; conjugate transpose}
\nt{$\fs$, $\fs'$, $\fs^{Y}_{D,i}$, $\fs^{W}_{D,i}$}{families of signs; those of \m{\spaceY{D}{i}} and \m{\spaceW{D}{i}}}
\nt{$G_\infty(\fs)$, $\mathrm U(\fs)$, $G_f(I)$, $K_0(I)$, $K_g(D)$}{the unitary groups at the infinite places and over the finite adeles; the integral points; the principal congruence subgroup at the chosen primes}
\nt{$g_v$, $k^{(v)}$, $\bar k$, $\tilde\gamma$, $x_{(v)}$}{component at $v$; the element with component $k$ at a split $v$; reduction; the chosen lift of \m{\gamma\in\sG(D)}; the element of a product with component $x$ at $v$}
\nt{$\cV=(J,\mathbf x)$, $\cV'$, $\cV''$, $\Gamma_\cV$}{spaces with a point; the rational points in the adelic group}
\nt{$F$, $F'$, $R(g)$, $\langle\ ,\ \rangle$}{forms; right translation; the Petersson product}
\nt{$\epsilon$, $\epsilon_l$, $\epsilon_{D,\fs}$, $u$, $\alpha$}{types; the kernel of a discrete series; the type of the data; an element of $G_\infty(\fs)$ or a vector; a function of compact support, and in $\alpha(D)$ the idele of the Weil data}
\nt{$l$, $l(\tau)$, $H$, $B$}{integer vectors (in Lemma~\ref{l:refl}, a vector over a finite field); a gap; an order of a cusp kernel (in Lemma~\ref{l:common}, with $B'$, a function on pairs of forms; in Lemma~\ref{l:Y}, a lattice basis)}
\nt{$\Sigma=(\varrho,\mathcal X,o)$, $\mathsf b_n$, $\mathsf b^*_n$, $z$}{a shape and its index set; $o$ also a Boolean, $z$ also a form or an idele}
\nt{\m{\besselIx{\Sigma}}, \m{\indexY{D}{i}{i<m(D)}}}{bijections of index sets (Definition~\ref{d:index})}
\nt{\m{\chosenIdele{D}}, \m{\ideleA{D}}, \m{\piIdele{D}}, \m{\latticeScale{D}{i}}}{the ideles of Definition~\ref{d:ideles}}
\nt{$\mathrm{Gr}(\beta,N)$, $\mathrm{Lb}(\beta,a,\pi,u)$, $s$}{graded Gram matrix and lattice basis; the length of the string}
\nt{\m{\spaceY{D}{i}}, \m{\spaceW{D}{i}}}{the spaces of Definition~\ref{d:YW}}
\nt{\m{\Prin{D}{\cV}}, \m{\prep{D}{\cV}}}{principal forms and the action of $\sG(D)$ (Definition~\ref{d:prin})}
\nt{$e_0(D)$, \m{\Good{D}{\cV}}, \m{\Row{D}}, \m{\rowDim{D}}}{the good idempotent, good forms, rows, their dimension}
\nt{$e_{I,c}$, $e^{\otimes}_{I,\mathbf c}$, $\mathrm{Good}_{M_\omega(D),r(D)}$, $x_\chi$}{idempotents of labels and of families, good tuples, the unit naming a character (note \emph{Primitive labels})}
\nt{$c$, $\mathbf c=(c_v)_v$, $c^{p,v}$, $\tilde e_p$, $X$, $\mathcal C$}{polynomials, families, cubic twists, the twisted idempotent; the variable; a set of families}
\nt{$\cH(D)$, $T$, $T_\cV$, $T^\gamma$, $\Pi$, $\Pi^p$, $\mathbb T^{v,n}$}{the Hecke algebra away from the chosen primes, its elements and their operators; the Hecke operators at a prime}
\nt{$p$, $p'$, $\mathcal F$, $\mathcal F'$, \m{\Prin{D}{\cV}[p]}}{eigensystems, finite sets of them, eigenparts}
\nt{$\Occ(D,\cV,p,v)$, \Rigid{D}{\cV}}{occurring primitive labels; Definition~\ref{d:rigid}}
\nt{\m{\levelPlaces{D}{\vartheta}}, \m{\levelMod{D}{\vartheta}{n}{o}}}{places and moduli of a tame character (Definition~\ref{d:levels})}
\nt{\m{\auxPair{D}{\vartheta}{i}}, $\kappa^{+}_{O}$, $\kappa^{-}_{O}$, $\phi_{D,i}$, $O$, $\nu_O$, $K_O$, $\psi_{O,w}$, $\rho_O$, $\rho_i$, $s_w$}{auxiliary pairs and their parts (Definition~\ref{d:aux}, Lemma~\ref{l:type})}
\nt{$\mathrm{Ral}_{D,\vartheta,i}$, $c_{D,\vartheta,i}(\gamma)$, $g_\gamma$, $L_{D,\vartheta}$, $A$, $A'$, $\Psi$, $\mathcal Z$}{Rallis terms and the objects of Lemma~\ref{l:common}}
\nt{$L$, $\Lambda_i$, $\Lambda_{\mathrm{ch}}$, $\Upsilon_w$, $\nu$, $\beta$, $n_1$, $\chi$, $\chi_n$, $\psi$}{lifts and the data of Input~6; $\beta$ also a bijection of index sets, $\chi$ a multiplicative character}
\nt{$\mathcal S$, $\pi$, $f$, $y$, $y_v$, $\rho$, $S$, $S'$}{a space of functions, a representation, a function, vectors or elements of group algebras, a representation, matrices}
\nt{$b_n$, $t$, $c_g$, $c^p_g$}{an orthonormal basis; the common trace of Input~1; coefficients of Hecke elements}
\nt{$i$, $i'$, $n$, $n'$, $a$, $k$, $k_1$, $k_2$, $h$, $\gamma$, $\iota$, $\iota'$, $\varpi$}{indices and natural numbers; an idele; integral matrices; a matrix over a finite field; elements of $\sG(D)$ or cosets; elements of index sets; a uniformizer}
\nt{$y_\tau$, $\mathsf T$}{frames at the infinite places; transpose}
\end{multicols}
\endgroup

\head{References}
{\fontsize{9.4}{12.6}\selectfont
All references are given as the library's files, the catalogue of this handout and its second reading record
them; I did not open the works.\par
J.~Arthur, L.~Clozel, \emph{Simple algebras, base change, and the advanced theory of the trace formula}, 1989;
Chapter~1, §2, Lemmas 2.4--2.6.\par
A.~Borel, \emph{Automorphic forms on reductive groups} (Park City lectures), §8.3 with §§7.6--7.9 (Godement's
lemma), for the finiteness of spaces of automorphic forms; at second hand.\par
R.~Jacobowitz, \emph{Hermitian forms over local fields}, Amer. J. Math. 84 (1962); the general classification.\par
J.-P.~Labesse, \emph{Changement de base CM et séries discrètes}, 2011; Théorèmes 5.1 and~5.4.\par
C.~T.~C.~Wall, \emph{On the classification of hermitian forms~I}, Compositio Math.~22 (1970); pp.~432--433 and
Theorem~2.\par}

\end{document}
